\chapter{Methodology}

\section{Problem formulation}
\label{sec:methodology_problem}

This section fixes the notation used throughout the chapter and the thesis, and
draws the distinction --- relied upon by every subsequent section --- between
mechanisms that act on training and mechanisms that act on inference.

\subsection{Taxonomy}

A \emph{taxonomy} is a rooted tree $\mathcal{T}$ whose nodes are partitioned into
$L$ levels, ordered from the most general to the most specific and indexed
$1\prec 2\prec\dots\prec L$. Level $l$ contains $C_l$ classes, and
$N:=\sum_{l=1}^{L}C_l$ denotes the total number of nodes. Two structural
assumptions are made throughout, and are verified once when the hierarchy graph
is built rather than at every training step:

\begin{enumerate}
    \item \textbf{Unique parent.} Every class at level $l+1$ has exactly one
    parent at level $l$. The parent map is written
    $\pi_l:\{1,\dots,C_{l+1}\}\to\{1,\dots,C_l\}$, and its matrix form is the
    incidence matrix $M_{l,l+1}\in\{0,1\}^{C_l\times C_{l+1}}$ with
    $M_{l,l+1}[i,j]=1$ if and only if $\pi_l(j)=i$. Uniqueness of the parent is
    exactly the statement that every column of $M_{l,l+1}$ contains a single
    non-zero entry.
    \item \textbf{No childless parent.} Every class at level $l$ has at least one
    child at level $l+1$, so that each row of $M_{l,l+1}$ is non-zero. As shown in
    Lemma~\ref{lemma:diag}, this is precisely the condition under which the HCC
    projection is well defined.
\end{enumerate}

Under these assumptions each leaf determines a unique root-to-leaf path, so a
label may equivalently be given as a leaf identifier or as the path
\[
y=(y_1,\dots,y_L),\qquad y_l\in\{1,\dots,C_l\},\qquad y_l=\pi_l(y_{l+1}).
\]
A path satisfying $y_l=\pi_l(y_{l+1})$ for all $l$ is called \emph{consistent};
the datasets provide consistent ground-truth paths by construction, whereas a
model's prediction need not be consistent, which is what the metrics of
\S\ref{ssec:inference_metrics} quantify.

All datasets used in this thesis are given a three-level hierarchy, $L=3$. The
methods are defined for arbitrary $L$ and the restriction is a property of the
experimental setting, not of the formulation; the two exceptions --- the HRN
baseline and the lexicographic mode of \S\ref{sec:methodology_lex}, both of which
require exactly three levels in their implementation --- are noted where they
arise.

\subsection{Classifier and score conventions}
\label{ssec:problem_scores}

A hierarchical classifier is a map $x\mapsto\big(z_1,\dots,z_L\big)$
parameterised by $\theta$, with
\[
z_l\;:=\;z_l(x;\theta)\in\mathbb{R}^{C_l}
\]
the vector of scores assigned to the classes of level $l$. The explicit
dependence on $x$ and $\theta$ is omitted whenever it is not needed. Levels are
always indexed coarse-to-fine, so that $z_1$ is the coarsest and $z_L$ the
finest.

Two conventions for producing $z_l$ occur among the model families compared here,
and the distinction matters when the constructions of the following sections are
applied.

\begin{itemize}
    \item \textbf{Per-level heads.} Most families (H-CAST, LH-DNN, HRN,
    HT-CapsNet) emit one score vector per level directly, as the output of a
    level-specific head applied to a shared representation.
    \item \textbf{Taxonomy-node scores.} Hier-COS instead emits a single node
    vector $u\in\mathbb{R}^{N}$ over the whole taxonomy and reduces it to level
    scores by subspace norms,
    \[
    z_l[c]=\big\lVert \Pi_{S_c}\,u\big\rVert_2,
    \]
    where $S_c$ is the taxonomy subspace assigned to node $c$ of level $l$. This
    reduction is non-linear, and consequently the affine constructions of
    \S\ref{sec:methodology_hcc} and \S\ref{sec:methodology_lhproj} act on the
    linear per-level node blocks of $u$ rather than on the subspace norms
    themselves. Where this distinction affects the interpretation of a result it
    is stated explicitly.
\end{itemize}

\subsection{Objective}
\label{ssec:problem_objective}

Unless a section states otherwise, training minimises a scalarised objective of
the form
\[
\mathcal{L}(\theta)
=
\sum_{l=1}^{L} w_l\,\ell_l\big(z_l,y_l\big)
\;+\;
\lambda\,\mathcal{R}\big(z_1,\dots,z_L\big),
\]
where $\ell_l$ is a level-wise loss (cross-entropy, or a family-specific
surrogate such as the capsule margin loss of HT-CapsNet), $w_l>0$ are level
weights with $\sum_l w_l=1$, and $\mathcal{R}$ is an optional family-specific
regulariser with coefficient $\lambda$.

The level weights are a genuine degree of freedom rather than an implementation
detail, and three settings are used: \texttt{equal} ($w_l=1/L$),
\texttt{kl\_leaf}, which concentrates weight on the finest level, and
\texttt{kl\_coarse}, its reversal. Because two of the mechanisms introduced below
either bypass $w_l$ or replace it by a uniform weighting, $w_l$ is treated as a
\emph{controlled variable} of the experimental design and is reported for every
run; the resulting confound between arms is discussed in
\S\ref{ssec:lex_setting} and in the threats to validity.

\subsection{Where a hierarchical method can act}
\label{ssec:problem_train_vs_inference}

A method for enforcing hierarchical structure may act at three distinct points of
the computation, and conflating them is the main source of imprecision in claims
about hierarchical classifiers. This distinction is the organising principle of
the chapter, and of the experimental design that closes it.

\begin{enumerate}
    \item \textbf{The forward map} $\theta\mapsto(z_1,\dots,z_L)$ may be modified,
    so that the scores themselves satisfy a structural property. Such an
    intervention is active at inference as well as during training, since it
    changes the function computed by a given parameter vector. This is the case of
    HCC (\S\ref{sec:methodology_hcc}).
    \item \textbf{The update direction} in parameter space may be modified,
    leaving the forward map untouched. Such an intervention changes the parameters
    the model converges to and has, by construction, no effect at inference. This
    is the case of lexicographic mode (\S\ref{sec:methodology_lex}) and, as shown
    in Proposition~\ref{prop:lhproj_forward}, of the LH-projection
    (\S\ref{sec:methodology_lhproj}).
    \item \textbf{The inference rule} that turns a trained network into a
    predicted path may be modified, leaving both the forward map and the training
    procedure untouched (\S\ref{sec:methodology_inference}). This point admits a
    further factorisation, which \S\ref{sec:methodology_inference} exploits: a
    \emph{readout}, which reduces the network's internal node coordinates to level
    scores, and a \emph{decoder}, which turns level scores into a path. Both can be
    replaced after training on a frozen checkpoint.
\end{enumerate}

The three loci are independent, and the experiments vary them one at a time so
that an observed effect can be attributed to the output contract, to the
optimisation trajectory, or to the inference rule.

The third locus deserves emphasis, because it supplies the instrument that makes
the other two interpretable. A mechanism that acts during training can improve the
reported metrics for either of two quite different reasons: it may have changed
the representation the network learned, or it may merely have supplied a better
rule for reading out a representation that was already there. Since a readout can
be swapped on a fixed checkpoint --- the same parameters, the same forward pass,
a different reduction to level scores --- the two explanations can be separated
by construction rather than argued about. This is the role of the post-hoc
inference grid of \S\ref{ssec:inference_grid}, and it is why the inference section
is not an evaluation appendix but a component of the argument.

\paragraph{Roadmap.} Sections~\ref{sec:methodology_hcc}
to~\ref{sec:methodology_lhproj} take the three priority mechanisms in turn, one
per locus: HCC on the forward map, lexicographic mode on the assembled update, the
LH-projection on the learning signal at a structural bottleneck. Each is derived,
its guarantee is stated precisely, and the scope of that guarantee is delimited;
the comparison between them is deferred, and made once, in
\S\ref{sec:methodology_design}. Section~\ref{sec:methodology_inference} then fixes
the inference rules, the metrics computed on top of them, and the
checkpoint-selection protocol. Section~\ref{sec:methodology_design} closes the
chapter by turning this apparatus into an experimental design: which combinations
are admissible, which research questions the admissible ones answer, which arms
exist solely to falsify an attribution, and which variables are held fixed for the
comparisons to mean what they claim.


\section{Hierarchical Constraint Cascade}
\label{sec:methodology_hcc}

The \emph{Hierarchical Constraint Cascade} (HCC) is a differentiable method for enforcing hierarchical relations among the outputs of a multi-level classifier. The method builds on the \textsc{HardNet} framework for imposing hard constraints on neural networks \cite{min2024hardnet}, of which it constitutes a specific instance: the equality-only case, applied recursively along a taxonomy.

This section first presents the general framework (\S\ref{ssec:hardnet}), then formally derives the hierarchical instantiation underlying HCC (\S\ref{ssec:hardnet_to_hcc}), relates the resulting construction to the lexicographic approaches that form the backbone of this thesis (\S\ref{ssec:lexicographic}), and finally discusses the extensions introduced here with respect to the original method (\S\ref{ssec:hcc_extensions}).

\subsection{The HardNet framework}
\label{ssec:hardnet}

\subsubsection{Motivation}

Prior knowledge about the input--output relations of a model is typically imposed in one of two ways. The first, \emph{soft constraints}, penalises violations through a regularisation term in the objective; this offers no guarantee of constraint satisfaction, particularly for inputs far from the training distribution, a condition that is critical in safety-critical applications \cite{min2024hardnet}. The second, \emph{hard constraints} imposed naively --- for instance by redefining the architecture so that the constraint holds by construction through an \emph{ad hoc} parameterisation --- risks a drastic reduction of the model's representational capacity.

HardNet was proposed as a resolution of this trade-off \cite{min2024hardnet}: a framework for constructing networks that satisfy \emph{input-dependent} hard constraints by construction, without sacrificing model capacity, and that remain trainable end-to-end with standard unconstrained optimisation algorithms.

\subsubsection{Definition}

Let $f_\theta:\mathcal{X}\to\mathbb{R}^{n_{\mathrm{out}}}$ be a neural network parameterised by $\theta$. HardNet enforces the input-dependent affine constraint

\[
b^l(x)\le A(x)f(x)\le b^u(x)\qquad\forall x\in\mathcal{X},
\]

with $A(x)\in\mathbb{R}^{n_c\times n_{\mathrm{out}}}$, by appending to $f_\theta$ a differentiable \emph{enforcement layer}, that is, a projection $\mathcal{P}$ such that $\mathcal{P}(f_\theta):\mathcal{X}\to\mathbb{R}^{n_{\mathrm{out}}}$ satisfies the constraint for every $x$ while remaining differentiable with respect to $\theta$ through $f_\theta$. The layer is agnostic to the upstream architecture: it can be attached to any network exposing an output vector of dimension $n_{\mathrm{out}}$ \cite{min2024hardnet}.

Under the assumptions that, for every $x$, the constraint is feasible and $A(x)$ has full row rank, the closed-form solution \textsc{HardNet-Aff} is given by \cite{min2024hardnet}

\[
\mathcal{P}(f_\theta)(x)
=
f_\theta(x)
+
A(x)^{+}
\Big[
\operatorname{ReLU}\!\big(b^l(x)-A(x)f_\theta(x)\big)
-
\operatorname{ReLU}\!\big(A(x)f_\theta(x)-b^u(x)\big)
\Big],
\]

where $A(x)^{+}:=A(x)^{\top}\big(A(x)A(x)^{\top}\big)^{-1}$ denotes the Moore--Penrose pseudoinverse. Two properties established in \cite{min2024hardnet} are directly relevant here:

\begin{itemize}
    \item \textbf{Exact constraint satisfaction}: for each row $a_i$ of $A(x)$, the projected output satisfies $a_i^{\top}\mathcal{P}(f_\theta)(x)=b^l_{(i)}(x)$ or $b^u_{(i)}(x)$ whenever the constraint is violated, and is left unchanged otherwise (Proposition 4.4 in \cite{min2024hardnet}).
    \item \textbf{Preservation of universal approximation}: if the function class realised by $f_\theta$ universally approximates a class $\mathcal{F}$, then the class of projected functions $\mathcal{P}(f_\theta)$ universally approximates the subset of $\mathcal{F}$ satisfying the constraint (Theorem 4.6 in \cite{min2024hardnet}). Enforcing the constraint by construction therefore does not, in the limit, reduce the representational capacity of the model.
\end{itemize}

\subsubsection{Rationale for the choice of HardNet}

Relative to the alternatives surveyed in \cite{min2024hardnet} --- soft constraints (no guarantee), iterative correction procedures such as DC3 (no guarantee within a fixed number of iterations, higher computational cost, sensitivity to the hyperparameters of the correction procedure) and parameterisations of the output set alone, which are input-independent --- HardNet-Aff is presented as the first method with a closed-form (hence inexpensive) forward pass capable of enforcing multiple \emph{input-dependent} affine constraints with a guarantee of exact satisfaction and without a proven loss of expressivity.

This combination of properties --- low computational cost, exact guarantee, architecture independence, and a proven preservation of model capacity --- is the reason this thesis adopts HardNet-Aff as the underlying mechanism for enforcing hierarchical consistency, rather than a soft penalty in the loss or an iterative correction procedure.

\subsection{Notation}
\label{ssec:formulation}

The notation of \S\ref{sec:methodology_problem} is used throughout: $z_l$ denotes
the level-$l$ logits of a classifier parameterised by $\theta$, levels are ordered
$1\prec 2\prec\dots\prec L$ from the most general to the most specific, and
$M_{l,l+1}\in\{0,1\}^{C_l\times C_{l+1}}$ is the incidence matrix of the
parent--child relation between consecutive levels, with exactly one non-zero entry
per column. Only the two structural assumptions of
\S\ref{sec:methodology_problem} --- unique parent, and no parent without children
--- are required by the derivation that follows.

\subsection{From the HardNet-Aff projection to the hierarchical cascade}
\label{ssec:hardnet_to_hcc}

HCC instantiates HardNet-Aff separately for each level transition $l\to l+1$, setting $A(x)\equiv M_{l,l+1}$ --- a \emph{fixed} matrix, not input-dependent, since it encodes the structure of the taxonomy rather than the input --- and imposing an equality-only constraint, that is $b^l=b^u=\hat z_l$:

\[
M_{l,l+1}\hat z_{l+1}=\hat z_l.
\]

An equality constraint is the degenerate case of the box constraint $b^l\le A z\le b^u$ obtained by setting $b^l=b^u$ (Remark 4.1 in \cite{min2024hardnet}). The following paragraph makes explicit the form taken by the general HardNet-Aff projection in this case.

It should be stressed that the constraint is imposed in \emph{logit space}. In general it does not imply the equality
\[
\operatorname{softmax}(\hat z_l)
=
M_{l,l+1}\operatorname{softmax}(\hat z_{l+1}).
\]
HCC must therefore not be read as a normalisation of probability mass between parent and children, but as a structural affine constraint on the representations fed to the objective function.

\subsubsection{Reduction to the equality-only case}

\begin{proposition}
For $b^l=b^u=b$, the HardNet-Aff projection reduces to
\[
\mathcal{P}(f_\theta)(x)=f_\theta(x)-A(x)^{+}\big(A(x)f_\theta(x)-b(x)\big).
\]
\end{proposition}
\begin{proof}
Let $v:=b-Af_\theta(x)$. For every scalar $v$ the identity $\operatorname{ReLU}(v)-\operatorname{ReLU}(-v)=v$ holds, as verified by cases: if $v\ge 0$ then $\operatorname{ReLU}(v)-\operatorname{ReLU}(-v)=v-0=v$; if $v<0$ then $=0-(-v)=v$. The identity extends componentwise to the vector case. Substituting into the general HardNet-Aff expression with $b^l=b^u=b$ yields $\operatorname{ReLU}(b-Af_\theta)-\operatorname{ReLU}(Af_\theta-b)=b-Af_\theta$, from which the claim follows.
\end{proof}

Applying this result with $A=M_{l,l+1}$, $f_\theta(x)=z_{l+1}$ and $b=\hat z_l$ recovers exactly the Euclidean projection onto the affine subspace $\{q:M_{l,l+1}q=\hat z_l\}$:

\[
\hat z_{l+1}
=
z_{l+1}
-
M_{l,l+1}^{\top}
\left(
M_{l,l+1}M_{l,l+1}^{\top}
\right)^{-1}
\left(
M_{l,l+1}z_{l+1}-\hat z_l
\right).
\]

By Proposition 4.4 in \cite{min2024hardnet}, specialised to $b^l=b^u$, this projection satisfies the constraint \emph{exactly} for any $z_{l+1}$, irrespective of whether the constraint was initially violated.

\subsubsection{A structural property specific to taxonomies}

Unlike the general case treated in \cite{min2024hardnet}, in which $A(x)$ may vary with the input and the invertibility of $A(x)A(x)^{\top}$ must be verified pointwise for every $x$, in HCC the matrix $M_{l,l+1}$ is fixed and its "single parent per child" structure admits a closed-form characterisation of $M_{l,l+1}M_{l,l+1}^{\top}$.

\begin{lemma}
\label{lemma:diag}
If every column of $M_{l,l+1}$ contains exactly one non-zero entry, then $M_{l,l+1}M_{l,l+1}^{\top}$ is diagonal, with $i$-th diagonal entry equal to the number of child classes of class $i$ at level $l$.
\end{lemma}
\begin{proof}
Write $M:=M_{l,l+1}$. For $i=i'$ we have $(MM^\top)_{i,i}=\sum_j M[i,j]^2=\sum_j M[i,j]$, since $M$ is binary, which equals the number $n_i$ of children of class $i$. For $i\ne i'$ we have $(MM^\top)_{i,i'}=\sum_j M[i,j]M[i',j]$; by hypothesis, for each column $j$ at most one of $M[i,j]$ and $M[i',j]$ is non-zero, so every summand vanishes and the entry is zero.
\end{proof}

It follows that $M_{l,l+1}M_{l,l+1}^{\top}=\operatorname{diag}(n_1,\dots,n_{C_l})$, which is invertible if and only if every class at level $l$ has at least one child at level $l+1$ --- a structural condition of the taxonomy, verifiable once when the hierarchy graph is built, rather than for every $x$ as required in the general HardNet-Aff setting. Moreover, since $M_{l,l+1}M_{l,l+1}^\top$ is diagonal, the numerically regularised term $M_{l,l+1}M_{l,l+1}^{\top}+\varepsilon I=\operatorname{diag}(n_1+\varepsilon,\dots,n_{C_l}+\varepsilon)$ remains diagonal: its inversion reduces to an elementwise operation that can be computed and cached before training, with no runtime matrix inversion. This is a computational advantage specific to the taxonomic case relative to the general variable-$A(x)$ setting in \cite{min2024hardnet}.

Note further that, for equality-only constraints, feasibility (Assumption 4.3-i in \cite{min2024hardnet}) is implied by full row rank (Assumption 4.3-ii): a surjective linear map admits a solution for any right-hand side. The only condition to be verified on the taxonomy is therefore the absence of parent classes without children.

\subsubsection{Recursive cascade}

Setting $\hat z_1=z_1$ and applying the projection recursively, with the stop-gradient operator $\operatorname{sg}(\cdot)$ on the right-hand side (discussed in \S\ref{ssec:hcc_extensions}), yields

\begin{align}
\hat z_1 &= z_1,\\
\hat z_2 &= P_{1,2}\big(z_2;\operatorname{sg}(\hat z_1)\big),\\
&\;\vdots\\
\hat z_L &= P_{L-1,L}\big(z_L;\operatorname{sg}(\hat z_{L-1})\big),
\end{align}

where each $P_{l,l+1}$ is, by construction, an instance of equality-only HardNet-Aff applied to the level pair $(l,l+1)$. The coarsest level is left unchanged and acts as the reference; each subsequent level is projected onto the affine set determined by the level above it, producing a coarse-to-fine cascade of hierarchical projections.

\subsection{Relation to lexicographic approaches}
\label{ssec:lexicographic}

Since this thesis is concerned with lexicographic formulations of hierarchical classification, the position of HCC within that framework requires explicit justification: HCC is not a lexicographic optimisation method, yet it realises --- in output space and in closed form --- the same priority structure that lexicographic formulations realise in objective space. This subsection makes the correspondence precise, together with its limits.

\subsubsection{Lexicographic priority as level-set restriction}

Given objectives ordered by priority $f_1\prec f_2\prec\dots\prec f_L$, lexicographic optimisation solves the sequence of problems

\[
\min_{q} f_k(q)
\qquad\text{subject to}\qquad
f_j(q)=f_j^{\star},\quad j<k,
\]

where $f_j^\star$ denotes the optimal value attained at stage $j$. The defining feature is not the ordering as such, but the fact that a lower-priority objective may act \emph{only within the level set already fixed by the higher-priority ones}. Equivalently: no finite weight can purchase a trade-off of a higher priority against a lower one, which is precisely what distinguishes lexicographic formulations from the weighted-sum scalarisations used by multi-level architectures such as B-CNN \cite{zhu2017bcnn}. This is the structure exploited by the lexicographic hybrid deep neural network of \cite{fiaschi2024informed}, which combines lexicographic multi-objective optimisation with non-standard analysis \cite{benci2022nonstandard,benci2019measure} to obtain a hierarchical classifier equipped with projection operators before each output layer.

\subsubsection{The structural correspondence}

The HCC projection derived in \S\ref{ssec:hardnet_to_hcc},

\[
\hat z_{l+1}
=
\arg\min_{q\in\mathbb{R}^{C_{l+1}}}
\tfrac{1}{2}\lVert q-z_{l+1}\rVert_2^2
\quad\text{subject to}\quad
M_{l,l+1}q=\hat z_l,
\]

has the same algebraic form. The coarse level contributes an equality constraint --- that is, a level set --- and the fine level is resolved, in the minimum-perturbation sense, strictly inside it. Three consequences follow, and they are the substantive sense in which HCC belongs in a thesis on lexicographic hierarchical classification:

\begin{itemize}
    \item \textbf{One-directional dominance.} The coarse level determines the feasible set for the fine level; the fine level cannot alter the coarse one through the projection branch. The stop-gradient operator extends this asymmetry from the forward pass to the gradient, so that the direction of influence is coarse-to-fine in both.
    \item \textbf{Non-negotiability.} Because the constraint is an equality enforced exactly (Proposition 4.4 in \cite{min2024hardnet}), coarse-level consistency is not traded against fine-level accuracy at any finite value of the level weights $w_l$. In this respect HCC behaves as a lexicographic priority rather than as a scalarised multi-objective compromise, even though the losses themselves are combined additively.
    \item \textbf{Closed form.} The restriction to the higher-priority level set, which in a lexicographic optimisation procedure is obtained by solving a constrained subproblem, is here available in closed form by Lemma~\ref{lemma:diag}, at the cost of a diagonal rescaling.
\end{itemize}

\subsubsection{Limits of the correspondence}

The analogy must not be overstated. Two differences are substantive and should be borne in mind when interpreting the experimental results.

First, lexicographic optimisation prioritises \emph{objectives} and requires the higher-priority problem to be solved to optimality before the next is considered. HCC constrains \emph{outputs}, and its anchor $\hat z_1=z_1$ is the network's current coarse prediction, not the coarse-optimal one; all levels continue to train jointly. HCC therefore exhibits lexicographic \emph{structure} without lexicographic \emph{sequential optimality}.

Second, the protection afforded by the stop-gradient is partial. It removes the influence of the fine-level loss on the coarse logits along the projection branch, but level-wise losses can still interact through the shared backbone, so parameter-space priority is not enforced.

\subsubsection{Methodological role in this thesis}

Given both the correspondence and its limits, HCC serves three distinct purposes within the present work.

\begin{enumerate}
    \item \textbf{It isolates the contribution of hierarchical priority from that of a particular lexicographic implementation.} The approach in \cite{fiaschi2024informed} couples lexicographic priority with a specific architecture and with machinery drawn from non-standard analysis. HCC delivers a comparable coarse-to-fine priority structure while remaining architecture-agnostic, allowing the two to be compared and the sources of any observed gain to be attributed.
    \item \textbf{It occupies the hard-constraint end of an axis whose soft end is lexicographic optimisation.} Lexicographic formulations realise priority at optimisation time, with guarantees that are asymptotic and schedule-dependent; HardNet realises it by construction, with exact satisfaction and a proven preservation of expressivity. Comparing the two therefore addresses whether hierarchical priority is better imposed on the optimisation trajectory or on the output contract.
    \item \textbf{It provides a controllable interpolation between the two regimes.} As described in \S\ref{ssec:hcc_extensions}, the activation coefficient $\alpha_t$ interpolates between unconstrained joint training and fully constrained training. A step schedule combined with freezing the coarse head at activation time would approximate a genuine two-stage lexicographic procedure; intermediate schedules relax it continuously. The schedule is thus not merely an optimisation device but the experimental knob along the axis described above.
\end{enumerate}

\subsection{Extensions of HCC with respect to HardNet}
\label{ssec:hcc_extensions}

The framework in \cite{min2024hardnet} is designed to impose a constraint --- or a set of constraints --- at a single output point of the network, typically specified by a known function of the problem (for instance a control barrier function, or the constraints of an optimisation problem). HCC adapts it to a different setting: a chain of $L-1$ projections applied in sequence, in which the right-hand side of each projection is not an externally known function but the \emph{projected output of the previous step of the same network}. This difference motivates two design choices that are not part of the original HardNet-Aff formulation, and which should not be attributed to its authors.

\paragraph{Stop-gradient.} In the original formulation the right-hand side $b(x)$ does not depend on $\theta$, so the question of gradient coupling between $b$ and the projection does not arise. In HCC, by contrast, $\hat z_l$ is itself the result of a projection that is differentiable with respect to $\theta$. Without further provision, the level-$(l+1)$ loss would backpropagate through the projection branch and directly modify $\hat z_l$, mixing the objectives of the two levels along that path. The operator $\operatorname{sg}(\cdot)$ --- which returns its argument unchanged in the forward pass and annihilates its gradient in the backward pass --- instead treats $\hat z_l$ as an anchor, realising the coarse-to-fine dependence discussed in \S\ref{ssec:lexicographic}. As noted there, this does not amount to full lexicographic protection of the shared parameters.

The stabilised projection actually used in the implementation is therefore

\[
\hat z_{l+1}
=
z_{l+1}
-
M_{l,l+1}^{\top}
\left(
M_{l,l+1}M_{l,l+1}^{\top}
+
\varepsilon I
\right)^{-1}
\left(
M_{l,l+1}z_{l+1}-\operatorname{sg}(\hat z_l)
\right).
\]

For $\varepsilon=0$ and full row rank the affine constraints are satisfied exactly; for $\varepsilon>0$ they are satisfied up to the numerical precision determined by the stabilisation term.

\paragraph{Activation curriculum.} Appendix C of \cite{min2024hardnet} documents that the enforcement layer can produce a vanishing gradient when all constraints active at a given point are simultaneously violated and the gradient of the loss with respect to the projected output lies in the span of the active constraint vectors --- a condition that becomes more likely as the number of constraints approaches the output dimension, as happens in the coarser and less populated levels of a taxonomy. As a mitigation, the same work proposes a simple \emph{warm start}: disabling the enforcement layer for the first $k$ epochs of training and activating it abruptly thereafter.

The coefficient $\alpha_t\in[0,1]$ introduced by HCC generalises this to a continuous transition. The logits actually used by the model are defined as

\[
z_l^{\mathrm{eff}}(t)
=
(1-\alpha_t)z_l
+
\alpha_t\hat z_l .
\]

A step schedule, with $\alpha_t\in\{0,1\}$, reproduces exactly the warm start of \cite{min2024hardnet}; linear, exponential or hyperbolic-tangent schedules realise a gradual version, with the further aim of attenuating the vanishing-gradient regime reported there. Since $\hat z_1=z_1$, the interpolation affects only the levels below the first. The coefficient determines three regimes:

\begin{itemize}
    \item for $\alpha_t=0$, HCC is inactive and the model coincides with the unconstrained classifier;
    \item for $0<\alpha_t<1$, the constraint is introduced gradually and is not yet satisfied exactly;
    \item for $\alpha_t=1$, only the projected logits are used and the affine constraint is fully active.
\end{itemize}

This formulation allows a comparison between applying HCC from the outset and deploying it as a delayed structural intervention.

\paragraph{Multi-level cascade.} Finally, whereas HardNet-Aff is applied and evaluated as a single terminal layer in \cite{min2024hardnet}, HCC applies it recursively across $L-1$ successive transitions. The guarantees of exact satisfaction (Proposition 4.4) and preservation of universal approximation (Theorem 4.6) established in \cite{min2024hardnet} concern a single application of the projection; their formal extension by induction over levels to the full cascade with $L>2$ is not proved in this work and remains an open theoretical question. Experimentally this means that exact satisfaction between two consecutive levels is guaranteed by each projection taken individually (Lemma~\ref{lemma:diag} and the preceding proposition), whereas the effect of composing the projections, the schedule $\alpha_t$ and the stop-gradient on the overall representational capacity must be assessed empirically.

\subsection{Objective function}

HCC does not necessarily introduce a new loss: it modifies the logits on which the objectives already defined by the classifier are evaluated. Writing $\ell_l$ for the loss associated with level $l$, the general objective is

\[
\mathcal{L}_{\mathrm{HCC}}
=
\sum_{l=1}^{L}
w_l\,\ell_l\big(z_l^{\mathrm{eff}},y_l\big)
+
\lambda\,\mathcal{R}\big(z_1^{\mathrm{eff}},\dots,z_L^{\mathrm{eff}}\big),
\]

where $w_l$ are optional level-wise weights and $\mathcal{R}$ an optional model-specific regularisation term. Because the projection is differentiable with respect to the logits, the model adapts during training to the corrected representations: before activation the loss is computed on the original logits, and after full activation on the projected ones.

\subsection{Architecture independence}

Like the HardNet enforcement layer, which is presented as applicable to any neural network provided its output contract is specified \cite{min2024hardnet}, HCC is defined on the output contract of the classifier rather than on its internal structure. Only the following are required:

\begin{enumerate}
    \item a vector of differentiable logits for each level of the hierarchy;
    \item a complete taxonomy, with a unique parent for every non-root class and no parent class without children (the condition required by Lemma~\ref{lemma:diag});
    \item a consistent ordering of the levels from coarsest to finest.
\end{enumerate}

The same cascade can therefore be attached to models with transformer, convolutional, capsule-based or taxonomic-frame backbones, provided their outputs can be reduced to this common interface. The formulation supports an arbitrary number of levels by recursive repetition of the same operation; a given implementation may restrict itself to a three-level hierarchy without that restriction belonging to the general definition of the method.

This independence concerns the applicability of the mechanism and does not imply that HCC produces the same effect on all architectures. The effectiveness of the constraint, as well as the interaction among cascade, curriculum and stop-gradient, must be verified experimentally for each combination of model, dataset, schedule and decoding strategy.


\section{Gradient-space lexicographic optimisation}
\label{sec:methodology_lex}

Section~\ref{sec:methodology_hcc} imposed hierarchical priority on the \emph{outputs} of the
classifier: the coarse logits define an affine set inside which the fine logits are resolved.
That construction has lexicographic structure, but, as noted in \S\ref{ssec:lexicographic}, it
leaves the shared parameters unprotected --- the level-wise losses still interact freely through
the backbone. This section describes the complementary mechanism implemented in this work, which
this thesis refers to as \emph{lexicographic mode}: priority is imposed directly on the
\emph{update direction} in parameter space, by projecting the gradients of lower-priority
objectives so that, to first order, they cannot act against the higher-priority ones.

The two mechanisms are configured independently (\texttt{hcc.enabled} and
\texttt{train.lexicographic.enabled}) and are varied one at a time in the experiments, so that
output-space and parameter-space priority can be attributed separately.

\subsection{Setting and requirements}
\label{ssec:lex_setting}

Lexicographic mode assumes a classifier exposing exactly three differentiable scalar level
objectives
\[
\ell_1(\theta),\quad \ell_2(\theta),\quad \ell_3(\theta),
\]
associated with the coarse, middle and fine levels and ordered by decreasing priority,
$\ell_1 \prec \ell_2 \prec \ell_3$. Writing
\[
g_l := \nabla_\theta \ell_l(\theta) \in \mathbb{R}^{P}
\]
for the gradient of the $l$-th objective with respect to the vector $\theta$ of trainable
parameters, each optimisation step computes the three gradients separately by three backward
passes through the shared graph, rather than a single backward pass on a scalarised total loss.

Two consequences of this design must be stated explicitly, because they affect the
interpretation of every experiment run in this mode.

\begin{itemize}
    \item \textbf{The scalarised loss is not used for the update.} In lexicographic mode the
    quantity actually applied to the parameters is a combination of the three per-level
    gradients; the weighted total loss is still computed and logged, but never backpropagated.
    Any term of the objective that is not decomposable into the three level losses therefore
    contributes nothing to the update. This is why the mode requires
    \texttt{model.loss.globalkl: false} for H-CAST: the global KL term is not part of any level
    objective and would be silently discarded. The analogous requirements for the baselines are
    that HRN uses its decomposed \texttt{level\_marginal} objective (with the leaf cross-entropy
    folded into the fine level), Hier-COS uses one of the decomposed
    \texttt{global\_softmax\_ce\_reg} / \texttt{level\_softmax\_ce\_reg} modes, and HT-CapsNet uses
    its three capsule margin losses. LH-DNN is excluded, since it does not expose three
    independent differentiable level objectives in the required form.
    \item \textbf{The level objectives entering the projection are the raw, unweighted ones.}
    The per-level weights $w_l$ of \S\ref{sec:methodology_hcc} scale the scalarised loss but not
    the tensors consumed by the projection. Activating lexicographic mode therefore also replaces
    the configured level weighting by a uniform one. This is a confound between the lexicographic
    and non-lexicographic arms and is treated as such in the threats to validity: where a
    baseline uses non-uniform weights, the comparison isolates ``priority plus uniform weighting''
    against ``weighted scalarisation'', not priority alone.
\end{itemize}

\subsection{Parameter blocks induced by gradient support}
\label{ssec:lex_blocks}

Projecting one gradient onto another is meaningful only on the coordinates where both objectives
actually act. In a hierarchical network the three losses do not reach the same parameters: a
level-specific head is reached by one objective only, while the backbone is reached by all three.
The implementation therefore partitions the parameters by their \emph{gradient support}
$S(i) := \{\,l : (g_l)_i \not\equiv 0\,\}$, where a coordinate is treated as outside the support of
$\ell_l$ when \texttt{autograd} reports it as unused for that objective:

\[
\begin{aligned}
T_1 &:= \{\, i : S(i) = \{1,2,3\} \,\} && \text{parameters reached by all three levels,}\\
T_2 &:= \{\, i : S(i) = \{1,2\} \,\}   && \text{reached by coarse and middle only,}\\
T_3 &:= \{\, i : S(i) = \{1\} \,\}     && \text{reached by the coarse level only,}\\
R   &:= \{\, i : \text{otherwise} \,\} && \text{remaining coordinates, typically level-specific heads.}
\end{aligned}
\]

For a block $B$ define the restricted inner product
$\langle u,v\rangle_B := \sum_{i\in B} u_i v_i$ and the induced norm
$\lVert u\rVert_B$. All projection coefficients below are computed with respect to these
restricted inner products, one block at a time, and coordinates in $R$ are left untouched:
where two objectives do not share a parameter there is no conflict to resolve, and rescaling a
head gradient by a coefficient estimated elsewhere in the network would be an artefact rather
than a correction. In the diagnostics of \texttt{docs/GRADIENT\_PARAM\_DIAGNOSTIC\_LOGS.md}
these blocks appear as \texttt{t1}, \texttt{t2}, \texttt{t3} and their unions
\texttt{t2t1}, \texttt{t3t2t1}.

\subsection{The projection operator}
\label{ssec:lex_projection}

Given a target gradient $g$, a reference gradient $r$ and a block $B$, the operator used
throughout is the Euclidean projection of $g$ onto the orthogonal complement of $r$, restricted
to $B$:

\[
\Pi_B(g;r)_i =
\begin{cases}
g_i - \dfrac{\langle g,r\rangle_B}{\langle r,r\rangle_B}\, r_i, & i \in B,\\[2ex]
g_i, & i \notin B.
\end{cases}
\]

The component along $r$ is removed unconditionally, irrespective of the sign of
$\langle g,r\rangle_B$, so that
\[
\langle \Pi_B(g;r), r\rangle_B = 0 .
\]
Lower-priority objectives are thus confined to the orthogonal complement of the higher-priority
direction, which is the parameter-space analogue of the level-set restriction of
\S\ref{ssec:lexicographic}: the subordinate objective may act only where the dominant one is, to
first order, indifferent. Removing the component also when the two gradients agree is what
distinguishes subordination from mere de-conflicting: an objective that is allowed to retain its
positive component along a higher-priority direction is still, in effect, being summed with it.

One guard completes the definition. If $\lVert r\rVert_B^2 \le \varepsilon$ the coefficient is
ill-conditioned and the projection is skipped, $\Pi_B(g;r) := g$; the default is
$\varepsilon = 10^{-12}$ (\texttt{train.lexicographic.eps}). Whether the guard fired is logged at
every step, so that a run can be audited for how often the nominal operator was actually applied.

\subsection{Composition modes}
\label{ssec:lex_modes}

With three objectives the pairwise operator can be composed in more than one way.
\texttt{train.lexicographic.projection\_mode} selects among three schemes; in all of them the
final update direction assigned to the optimiser is the sum
$d := \tilde g_1 + \tilde g_2 + \tilde g_3$ of the (possibly projected) level gradients.

\paragraph{\texttt{coarse\_first} (default).} The coarse gradient is left unchanged; the middle
gradient is orthogonalised against it, separately on $T_2$ and on $T_1$; the fine gradient is
orthogonalised, on $T_1$, against the \emph{resultant} of the two higher-priority directions:
\[
\tilde g_1 := g_1,\qquad
\tilde g_2 := \Pi_{T_2}\!\big(g_2;g_1\big)\ \text{on }T_2,\quad
m := \Pi_{T_1}\!\big(g_2;g_1\big)\ \text{on }T_1,\qquad
\tilde g_3 := \Pi_{T_1}\!\big(g_3;\,g_1+m\big).
\]

\paragraph{\texttt{pairwise\_orthogonal}.} Identical for the middle level, but the fine gradient
is orthogonalised sequentially against each higher-priority direction rather than against their
resultant:
\[
\tilde g_3 := \Pi_{T_1}\!\Big(\Pi_{T_1}\big(g_3;g_1\big);\, m\Big).
\]

\paragraph{\texttt{fine\_first}.} The priority order is inverted: the fine gradient is left
unchanged, the middle gradient is orthogonalised against it, and the coarse gradient is
orthogonalised against the resulting higher-priority directions. This mode is not a hierarchical
hypothesis but a \emph{control}: it applies exactly the same amount of gradient surgery with the
opposite ordering, and so separates the effect of hierarchical priority from the effect of
orthogonalisation as such. Any gain observed under \texttt{coarse\_first} that is reproduced
under \texttt{fine\_first} cannot be attributed to the hierarchy.

The distinction between the first two modes is not cosmetic, and the following proposition makes
it precise.

\begin{proposition}
\label{prop:lex_orthogonality}
Fix the block $T_1$ and assume all degeneracy guards inactive. Then, under
\texttt{coarse\_first},
\[
\langle \tilde g_3,\; g_1+m\rangle_{T_1}=0,
\]
whereas under \texttt{pairwise\_orthogonal} the stronger conclusion
\[
\langle \tilde g_3, g_1\rangle_{T_1}=\langle \tilde g_3, m\rangle_{T_1}=0
\]
holds, that is, $\tilde g_3 \perp \operatorname{span}\{g_1,m\}$ on $T_1$.
\end{proposition}
\begin{proof}
The first claim is the defining property of $\Pi_{T_1}(\cdot\,;g_1+m)$. For the second, write
$f' := \Pi_{T_1}(g_3;g_1)$, so $\langle f',g_1\rangle_{T_1}=0$, and
$\tilde g_3 = f' - c\,m$ with $c=\langle f',m\rangle_{T_1}/\lVert m\rVert^2_{T_1}$, so
$\langle \tilde g_3,m\rangle_{T_1}=0$. Since $m=\Pi_{T_1}(g_2;g_1)$ satisfies
$\langle m,g_1\rangle_{T_1}=0$, it follows that
$\langle \tilde g_3,g_1\rangle_{T_1}=\langle f',g_1\rangle_{T_1}-c\langle m,g_1\rangle_{T_1}=0$.
\end{proof}

Removing the component along a resultant is therefore weaker than removing the components along
the two directions that generate it: under \texttt{coarse\_first} the fine gradient may still
possess a component along $g_1$, compensated by an opposite component along $m$. The logged
quantities \texttt{post\_cos\_t1\_fine\_proj\_coarse} and
\texttt{post\_cos\_t1\_fine\_proj\_mid\_proj} measure exactly this residual, and are expected to
vanish only in \texttt{pairwise\_orthogonal}.

\subsection{First-order priority guarantee}
\label{ssec:lex_guarantee}

The sense in which this construction realises lexicographic priority is a first-order,
per-step one, and it is worth stating exactly.

\begin{proposition}
\label{prop:lex_firstorder}
Fix the block $T_1$, assume \texttt{pairwise\_orthogonal} with all degeneracy guards inactive,
and let $d=g_1+m+\tilde g_3$ be the assigned direction. Then
\[
\text{(i)}\quad \langle d, g_1\rangle_{T_1}=\lVert g_1\rVert^2_{T_1},
\qquad
\text{(ii)}\quad \langle d, g_2\rangle_{T_1}=\langle g_1,g_2\rangle_{T_1}+\lVert m\rVert^2_{T_1}.
\]
\end{proposition}
\begin{proof}
(i) By Proposition~\ref{prop:lex_orthogonality}, $\langle \tilde g_3,g_1\rangle_{T_1}=0$, and by
construction $\langle m,g_1\rangle_{T_1}=0$. (ii) Write $g_2=m+c\,g_1$ with
$c=\langle g_2,g_1\rangle_{T_1}/\lVert g_1\rVert^2_{T_1}$. Then
$\langle m,g_2\rangle_{T_1}=\lVert m\rVert^2_{T_1}$ and, again by
Proposition~\ref{prop:lex_orthogonality},
$\langle \tilde g_3,g_2\rangle_{T_1}=\langle \tilde g_3,m\rangle_{T_1}+c\langle \tilde
g_3,g_1\rangle_{T_1}=0$; adding $\langle g_1,g_2\rangle_{T_1}$ gives the claim.
\end{proof}

Since the directional derivative of $\ell_l$ along the step $\theta\mapsto\theta-\eta d$ is
$-\eta\langle g_l,d\rangle+O(\eta^2)$, the proposition reads as follows.

\begin{itemize}
    \item By (i), the first-order decrease of the coarse objective on the shared trunk is exactly
    the one its own gradient would produce: neither the middle nor the fine objective can slow it,
    whatever their magnitude. No finite rescaling of the lower levels can purchase a trade-off
    against the coarse level --- the defining property of a lexicographic, as opposed to
    scalarised, formulation.
    \item By (ii), the middle objective receives a guaranteed non-negative contribution
    $\lVert m\rVert^2_{T_1}$ from its own projected gradient plus a term
    $\langle g_1,g_2\rangle_{T_1}$ of arbitrary sign contributed by the coarse level alone. The
    fine level contributes nothing. Degradation of a level by a \emph{higher}-priority one is
    permitted; degradation by a \emph{lower}-priority one is excluded. This asymmetry is precisely
    the priority structure that \S\ref{ssec:lexicographic} identified as the substance of the
    lexicographic ordering.
    \item The same two identities hold verbatim on $T_2$ for the pair $(\ell_1,\ell_2)$, with
    $\tilde g_3$ absent; on $T_3$ the direction is $g_1$ unmodified; on $R$ each coordinate keeps
    its own level gradient.
\end{itemize}

Four limitations bound the scope of this guarantee and should be carried into the interpretation
of the results.

\begin{enumerate}
    \item \textbf{First-order and per-step.} The statement concerns directional derivatives at the
    current $\theta$, not the value of the objectives after the step: it guarantees a descent
    \emph{direction} for the dominant objective, not monotone decrease, and says nothing about the
    limit of the trajectory. It is therefore not sequential lexicographic optimality: the coarse
    objective is never solved to optimality before the others are considered. As with HCC, this is
    lexicographic \emph{structure} without lexicographic \emph{sequential optimality}.
    \item \textbf{The optimiser is not gradient descent.} The projected direction $d$ is written
    into \texttt{param.grad} and consumed by AdamW, whose momentum and per-coordinate
    preconditioning do not preserve inner products. The applied update is therefore not $-\eta d$,
    and Proposition~\ref{prop:lex_firstorder} constrains the assigned gradient rather than the
    realised step. Weight decay acts on top of both.
    \item \textbf{Stochasticity.} All gradients are mini-batch estimates. Orthogonality is enforced
    between the estimates, not between the population gradients, and the projection coefficients
    inherit their variance.
    \item \textbf{Block restriction.} The guarantee holds inside each block. It says nothing about
    coordinates in $R$, where the level objectives do not compete, nor about interactions mediated
    by the forward pass rather than by the shared parameters.
\end{enumerate}

\subsection{Activation epoch}
\label{ssec:lex_activation}

As with the HCC curriculum of \S\ref{ssec:hcc_extensions}, the mechanism can be delayed:
\texttt{train.lexicographic.start\_epoch} defines the first epoch at which the projection is
active. Before it, training proceeds by ordinary backpropagation of the scalarised weighted loss;
from it onwards, the update is the projected combination above. The transition is a step, not a
ramp: unlike $\alpha_t$, there is no continuous interpolation between the two regimes, so a run
with \texttt{start\_epoch} $=k$ is exactly the baseline for its first $k$ epochs. This makes the
comparison between \texttt{step@0} and a late activation such as \texttt{step@80} a comparison
between priority imposed from initialisation and priority imposed as a late intervention on an
already-trained representation, with the two arms sharing, in expectation, their entire early
trajectory.

\subsection{Implementation notes}
\label{ssec:lex_implementation}

\paragraph{Cost.} Each optimisation step performs three backward passes over the shared graph
instead of one, and holds three gradient copies of the trainable parameters simultaneously. The
same three per-level gradients also feed the diagnostics, so they are computed once and reused
rather than recomputed.

\paragraph{Mixed precision.} Under automatic mixed precision the loss is scaled before the
backward pass, which would multiply every gradient by the current scale factor. Since a
projection coefficient is a ratio of inner products, it is invariant to a common positive factor;
the implementation nevertheless computes the projection in unscaled units --- so that the logged
coefficients, cosines and norms are directly comparable across steps with different scale factors
--- and multiplies the resulting direction by the current scale before assigning it, so that the
gradient scaler can unscale it and perform its usual non-finite check.

\paragraph{Diagnostics.} Because the guarantee above is conditional on the degeneracy guards
being inactive and on the mode selected, the implementation logs, per epoch: the pre-projection
cosines between level gradients on each block (\texttt{cos\_t1\_mid\_coarse},
\texttt{cos\_t1\_fine\_higher}, \dots), which measure how much conflict there is to resolve; the
post-projection cosines, which verify that the intended orthogonality was achieved; per-block
gradient norms before and after projection, which quantify how much of the lower-priority signal
was removed; binary flags recording whether each projection step was applied or skipped; and
per-block parameter norms and epoch-to-epoch parameter displacements. The full key glossary is
given in \texttt{docs/GRADIENT\_PARAM\_DIAGNOSTIC\_LOGS.md}. Claims about lexicographic behaviour
in the experimental chapters are grounded in these logged quantities rather than in configuration
flags.

\subsection{Position relative to HCC}
\label{ssec:lex_vs_hcc}

Lexicographic mode and HCC realise the same coarse-to-fine ordering at two different loci of
\S\ref{ssec:problem_train_vs_inference}. HCC constrains the logits: the restriction is exact,
holds per sample, and holds at inference as well as during training, but it governs only what
the heads emit and leaves the shared parameters free to be pulled in incompatible directions by
the level losses. Lexicographic mode constrains the assembled update: the restriction is
first-order and holds per batch, it says nothing about the scores a given parameter vector
produces, and it has no effect at inference whatsoever.

Whether the two can be combined is a question about the objective, not about the loci, and it is
settled in \S\ref{ssec:design_axis}: activating HCC makes each level loss a function of every
level's raw logits, which destroys precisely the separability into three level objectives that
the construction of \S\ref{ssec:lex_setting} presupposes. The full comparison of the three
mechanisms, including the LH-projection of \S\ref{sec:methodology_lhproj}, is given once in
Table~\ref{tab:framework_mechanisms}.


\section{LH-projection}
\label{sec:methodology_lhproj}

The third mechanism considered in this thesis originates in the lexicographic
hybrid deep neural network of \cite{fiaschi2024informed}, and is referred to here
as the \emph{LH-projection}. Where HCC constrains the logits emitted by the heads
(\S\ref{sec:methodology_hcc}) and lexicographic mode constrains the update
direction assembled from three separate backward passes
(\S\ref{sec:methodology_lex}), the LH-projection operates at the branch points of
a shared trunk: before each level head reads the shared representation, the
component of that representation lying in the subspace already claimed by the
higher-priority heads is removed.

This section derives the operator and its properties
(\S\ref{ssec:lhproj_construction}--\S\ref{ssec:lhproj_forward}), states the
instantiation used here on top of the Hier-COS taxonomy-frame architecture,
including the structural conditions that instantiation requires
(\S\ref{ssec:lhproj_hiercos}), and delimits the guarantee
(\S\ref{ssec:lhproj_limits}).

\subsection{The branch projection}
\label{ssec:lhproj_construction}

Consider a classifier with a shared trunk producing a pre-activation
$k\in\mathbb{R}^{D}$ and a representation $z=\rho(k)\in\mathbb{R}^{D}$, followed by
$L$ linear level heads with weight matrices $W_l\in\mathbb{R}^{C_l\times D}$. In
the unconstrained model every head reads the same $z$, so the level objectives
compete for the same directions of $\mathbb{R}^{D}$ with no ordering between them.

The LH-projection introduces an ordering by treating the rows of the
higher-priority heads as \emph{reserved directions}. For level $l>1$ define the
stacked matrix
\[
A^{(l)}
:=
\begin{bmatrix} W_1\\ \vdots\\ W_{l-1}\end{bmatrix}
\in\mathbb{R}^{R_l\times D},
\qquad
R_l:=\sum_{j<l}C_j ,
\]
whose rows are the class vectors of every level above $l$, and let
\[
c^{(l)}
:=
\big(A^{(l)}\big)^{\!\top}
\Big(A^{(l)}\big(A^{(l)}\big)^{\!\top}+\varepsilon I\Big)^{-1}
A^{(l)} z
\;=:\;
P^{(l)}_\varepsilon\, z .
\]
The representation offered to head $l$ is $z-c^{(l)}$: the higher levels keep the
subspace they span, and level $l$ is confined to what remains. The rows of
$A^{(l)}$ are taken \emph{detached} from the computational graph, so that the
reserved subspace is treated as given rather than as something the lower level may
reshape --- the parameter-space counterpart of the stop-gradient of
\S\ref{ssec:hcc_extensions}.

The regularisation $\varepsilon I$ plays the same role as in
\S\ref{ssec:hcc_extensions}: it makes the linear system solvable when the stacked
head rows are close to linearly dependent, at the price of an exactness that the
following proposition quantifies.

\begin{proposition}
\label{prop:lhproj_operator}
Let $A\in\mathbb{R}^{R\times D}$ have singular values $\sigma_1,\dots,\sigma_r>0$
with right singular vectors $v_1,\dots,v_r$, and let
$P_\varepsilon:=A^{\top}(AA^{\top}+\varepsilon I)^{-1}A$ with $\varepsilon>0$.
Then $P_\varepsilon$ is symmetric positive semi-definite with
\[
P_\varepsilon=\sum_{i=1}^{r}\frac{\sigma_i^2}{\sigma_i^2+\varepsilon}\,v_iv_i^{\top},
\]
so that its eigenvalues are $\sigma_i^2/(\sigma_i^2+\varepsilon)\in(0,1)$ on
$\operatorname{row}(A)$ and $0$ on $\operatorname{row}(A)^{\perp}$. Consequently
$P_\varepsilon\to P_0$, the orthogonal projector onto $\operatorname{row}(A)$, as
$\varepsilon\to 0^{+}$, and for $\varepsilon>0$ the residual retained by
$I-P_\varepsilon$ along $v_i$ is $\varepsilon/(\sigma_i^2+\varepsilon)$.
\end{proposition}
\begin{proof}
Write the thin SVD $A=U\Sigma V^{\top}$ with $U^{\top}U=V^{\top}V=I_r$ and
$\Sigma=\operatorname{diag}(\sigma_1,\dots,\sigma_r)$. Then
$AA^{\top}=U\Sigma^{2}U^{\top}$ and
\[
(AA^{\top}+\varepsilon I)^{-1}
=U(\Sigma^{2}+\varepsilon I)^{-1}U^{\top}+\tfrac{1}{\varepsilon}\big(I-UU^{\top}\big).
\]
Since $(I-UU^{\top})U=0$, the second term is annihilated on both sides by
$A^{\top}=V\Sigma U^{\top}$ and $A=U\Sigma V^{\top}$, leaving
$P_\varepsilon=V\Sigma(\Sigma^{2}+\varepsilon I)^{-1}\Sigma V^{\top}
=V\operatorname{diag}\big(\sigma_i^2/(\sigma_i^2+\varepsilon)\big)V^{\top}$,
which is the stated spectral decomposition. The remaining claims follow by
inspection of the eigenvalues.
\end{proof}

Two consequences are used later. First, the removal is exact only in the limit
$\varepsilon\to0$: for $\varepsilon>0$ a fraction
$\varepsilon/(\sigma_i^2+\varepsilon)$ of the reserved component survives, and the
fraction is largest along the \emph{weakest} singular directions of the stacked
head matrix --- that is, precisely where the higher-priority heads are least
confident of their geometry. Second, $I-P_\varepsilon$ is never rank-deficient for
$\varepsilon>0$, so the operator is always well defined; whether it is
\emph{informative} is a separate question, settled by the rank condition of
\S\ref{ssec:lhproj_hiercos}.

\paragraph{Sample-dependent variant.} In the formulation of
\cite{fiaschi2024informed} the reserved subspace is not the row space of the raw
head weights but that of the weights composed with the derivative of the shared
activation at the current input,
\[
A^{(l)}[b]:=\big[W_1;\dots;W_{l-1}\big]\operatorname{diag}\!\big(\rho'(k_b)\big),
\]
evaluated per sample $b$ of the batch. The motivation is that a head acts on the
trunk through $\rho$, so the directions it actually controls at a given input are
those surviving the activation's local Jacobian. This makes the protected subspace
input-dependent, and the linear system must be solved per sample rather than once
per batch. Both variants are implemented and compared:
\texttt{model.projection.rho\_enabled} selects the sample-dependent form, whose
cost is a batch of $R_l\times R_l$ solves per level, against the batch-shared form
obtained by omitting $\rho'$. In the LH-DNN reference implementation $\rho$ is a
ReLU and $\rho'(k)=\mathbb{1}[k>0]$ is a binary mask; in the Hier-COS
instantiation of \S\ref{ssec:lhproj_hiercos} a shared PReLU is used and
$\rho'(k)=\mathbb{1}[k>0]+a\,\mathbb{1}[k\le 0]$ with $a$ the learned negative
slope, taken detached.

\subsection{What the projection acts on}
\label{ssec:lhproj_forward}

The implementation does not feed $z-c^{(l)}$ to head $l$. It feeds
\[
z^{(l)} \;:=\; z-c^{(l)}+\operatorname{sg}\big(c^{(l)}\big),
\]
and this changes the character of the mechanism in a way that must be made
explicit, because it determines where the LH-projection belongs among the three
interventions of \S\ref{ssec:problem_train_vs_inference}.

\begin{proposition}
\label{prop:lhproj_forward}
Let $c^{(l)}=P^{(l)}_\varepsilon z$ with $P^{(l)}_\varepsilon$ constant with
respect to $z$, and let
$z^{(l)}=z-c^{(l)}+\operatorname{sg}(c^{(l)})$. Then
\[
\text{(i)}\quad z^{(l)}=z
\qquad\text{(forward value)},
\qquad\qquad
\text{(ii)}\quad
\frac{\partial z^{(l)}}{\partial z}=I-P^{(l)}_\varepsilon
\qquad\text{(Jacobian)} .
\]
\end{proposition}
\begin{proof}
(i) The stop-gradient operator returns its argument unchanged in the forward pass,
so $z^{(l)}=z-c^{(l)}+c^{(l)}=z$. (ii) In the backward pass
$\operatorname{sg}(c^{(l)})$ contributes no derivative, so
$\partial z^{(l)}/\partial z=I-\partial c^{(l)}/\partial z$. The matrix
$P^{(l)}_\varepsilon$ is built from detached head weights and, in the
sample-dependent variant, a detached activation derivative; it is therefore
constant with respect to $z$, and $c^{(l)}=P^{(l)}_\varepsilon z$ is linear in
$z$ with $\partial c^{(l)}/\partial z=P^{(l)}_\varepsilon$.
\end{proof}

The mechanism is therefore \emph{not} a modification of the forward map. Head $l$
receives the unprojected representation $z$, and the function computed by a given
parameter vector is exactly that of the unconstrained multi-head classifier. What
the projection modifies is the gradient: by (ii), the contribution of $\ell_l$ to
the gradient reaching the shared trunk is
\[
\big(I-P^{(l)}_\varepsilon\big)\,\nabla_{z}\,\ell_l ,
\]
that is, the component of the level-$l$ learning signal lying along the class
directions already claimed by the levels above it is removed before it reaches the
shared parameters. In the terms of \S\ref{ssec:problem_train_vs_inference}, the
LH-projection is an intervention on the update direction, not on the forward map,
and it has no effect at inference.

This places it beside lexicographic mode rather than beside HCC, and the two admit
a direct comparison, which is one reason both are studied here:

\begin{itemize}
    \item \textbf{Reference directions.} Lexicographic mode orthogonalises against
    the \emph{gradients} of the higher-priority objectives, which change at every
    step and depend on the current batch. The LH-projection orthogonalises against
    the \emph{head weight rows}, a slowly varying geometric object that encodes
    where the higher levels have decided their classes live. The former asks
    ``in which direction does the coarse loss want to move?''; the latter asks
    ``which part of the representation does the coarse level own?''.
    \item \textbf{Granularity.} The LH-projection is applied per sample, inside the
    forward graph, before any gradient is accumulated; lexicographic mode is
    applied per batch, to the accumulated gradient, on parameter blocks determined
    by gradient support (\S\ref{ssec:lex_blocks}).
    \item \textbf{Locus.} The LH-projection filters the signal at a single point,
    the branch point of the shared trunk, and therefore cannot control how the
    level objectives interact in parameters upstream of it beyond what that filter
    removes. Lexicographic mode acts directly on the parameters of every shared
    block.
    \item \textbf{Cost.} The LH-projection requires one backward pass and $L-1$
    linear solves of size $R_l$ (per sample in the $\rho'$ variant); lexicographic
    mode requires $L$ backward passes and no solves.
\end{itemize}

\subsection{Instantiation on Hier-COS}
\label{ssec:lhproj_hiercos}

The LH-projection is defined for a shared trunk read by independent linear level
heads. Hier-COS as published does not have that structure: it emits a single node
vector, maps it through a \emph{fixed} frame shared by the whole taxonomy, and
scores levels by subspace norms (\S\ref{ssec:problem_scores}). Transplanting the
projection onto it therefore requires three structural changes, which are part of
the construction rather than free hyperparameters, and which are stated here for
that reason.

\paragraph{(i) Independent level heads, hence an identity frame.} The projection
is defined in terms of the weight matrices $W_1,\dots,W_{l-1}$ of the heads above
level $l$. Under the dense random orthonormal frame of the published model the
scores of every level are linear mixtures of a single shared node space, and no
per-level matrix exists whose rows could be stacked into $A^{(l)}$: the object the
construction needs is simply not present. The projected variant therefore replaces
the frame by the identity --- equivalently, by a per-level block-diagonal frame,
which leaves the level blocks independent --- and attaches one learnable linear
head per level to the transform output. The implementation enforces this: enabling
the projection with a dense frame is rejected at construction. The frame is thus a
\emph{precondition} of the method, and the identity-versus-random comparison
measures the cost of the precondition, not a tuning choice; that cost is reported
in the experimental chapter.

\paragraph{(ii) Decomposed level objectives.} For the same reason, the objective
must expose one differentiable term per level. The published global objective
applies a single softmax over all $N$ taxonomy nodes and yields one scalar, from
which no per-level signal can be extracted at the branch point. The projected
variant therefore uses the level-wise objective, in which the softmax is taken
within each level's node block, together with the same node-norm regulariser. As
with the frame, the global-versus-level comparison quantifies the price of a
structural requirement.

\paragraph{(iii) A non-trivial complement.} The projection is informative only if
the reserved subspace does not exhaust the representation.

\begin{proposition}
\label{prop:lhproj_rank}
Let $D$ be the width of the shared representation. If $D\le R_L=\sum_{l<L}C_l$
then, for $A^{(L)}$ of full row rank, $\operatorname{row}(A^{(L)})=\mathbb{R}^{D}$
and $I-P^{(L)}_0=0$: the projection annihilates the entire learning signal of the
finest level. A non-trivial complement at every level requires
$D>\sum_{l<L}C_l$.
\end{proposition}
\begin{proof}
$A^{(L)}$ has $R_L$ rows in $\mathbb{R}^{D}$; if it has full row rank and
$R_L\ge D$ its row space is all of $\mathbb{R}^{D}$, so the orthogonal projector
onto that row space is the identity and its complement is zero. Since $R_l$ is
increasing in $l$, the binding constraint is $l=L$.
\end{proof}

The condition $D>\sum_{l<L}C_l$ is checked at construction and controls
\texttt{model.projection.feature\_dim}, which decouples the width of the
representation from the taxonomy size $N$ used by the unprojected model. It is
worth noting that for $\varepsilon>0$ the failure predicted by
Proposition~\ref{prop:lhproj_rank} is not a hard zero but a uniform shrinkage of
the fine-level gradient by $\varepsilon/(\sigma_i^2+\varepsilon)$
(Proposition~\ref{prop:lhproj_operator}), which would be hard to distinguish from
a badly scaled learning rate; the explicit check is preferred to relying on the
diagnostic.

\paragraph{Remaining degrees of freedom.} Given the three requirements above, the
following are genuine ablation axes, varied in the experimental chapter and not
constitutive of the method: the placement of the projection relative to the
transformation layer (\texttt{model.transform\_mode}, i.e.\ whether the projected
representation is the full residual transform, a batch-norm-plus-linear
reduction, or the backbone output itself); the sample-dependent versus
batch-shared reserved subspace (\texttt{model.projection.rho\_enabled}); the
representation width, subject to Proposition~\ref{prop:lhproj_rank}
(\texttt{model.projection.feature\_dim}); and the \emph{advantage}
parameterisation, in which the logits of level $l$ are expressed as a residual on
top of the detached logit of the predicted parent,
$z_l \leftarrow z_l + z_{l-1}[\pi_{l-1}(\cdot)]$, so that the head predicts a
correction to its parent rather than an absolute score
(\texttt{model.projection.advantage\_enabled}).

\paragraph{Exclusivity with HCC.} The LH-projection and HCC are mutually exclusive
on Hier-COS, and enabling both is rejected at construction. The reason is that they
would act on the same quantity in incompatible ways: the LH-projection determines
the per-level logits by filtering the representation each head reads, while HCC
overwrites those same logits with the solution of an affine problem anchored on the
level above. Composing them would make the reserved subspace of the LH-projection
--- built from the head weights --- describe a map that no longer determines the
logits actually used by the objective.

\subsection{Scope of the guarantee}
\label{ssec:lhproj_limits}

The guarantee attaching to the LH-projection in \cite{fiaschi2024informed} is
narrower than the construction may suggest, and four qualifications should be
carried into the interpretation of the results. They parallel those of
\S\ref{ssec:lex_guarantee}.

\begin{enumerate}
    \item \textbf{First-order.} The statement concerns the gradient at the current
    parameters. It guarantees that the level-$l$ signal reaching the trunk carries
    no component along the reserved subspace at that point; it does not guarantee
    that the higher-priority objectives do not increase after the step, and says
    nothing about the limit of the trajectory. As with HCC and with lexicographic
    mode, this is lexicographic \emph{structure} without lexicographic
    \emph{sequential optimality}.
    \item \textbf{Per-sample and branch-point-local.} The projection is applied at
    the branch point, on the representation each head reads. Interactions between
    level objectives that are mediated by parameters upstream of the branch point,
    or by batch statistics such as those of the normalisation layers, are outside
    its scope.
    \item \textbf{Conditional on the activation.} The sample-dependent form
    presumes that composing the head weights with $\rho'$ correctly identifies the
    directions the higher levels control, which holds for a piecewise-linear
    $\rho$ whose derivative is locally constant. It is exact for ReLU, where
    $\rho'$ is a binary mask, and approximate for the PReLU used here, whose
    negative slope is itself learned and is taken detached in the projection.
    \item \textbf{Approximate, and least accurate where it matters most.} By
    Proposition~\ref{prop:lhproj_operator} the removal leaves a residual
    $\varepsilon/(\sigma_i^2+\varepsilon)$ along the $i$-th singular direction of
    the stacked head matrix. The default $\varepsilon=10^{-6}$ makes this
    negligible for well-conditioned heads, but the residual grows as $\sigma_i\to
    0$, that is, early in training and along class directions the higher levels
    have not yet separated.
\end{enumerate}


\section{Hierarchical inference strategies}
\label{sec:methodology_inference}

The preceding sections described how a network is trained. Turning the trained
network into a predicted path is a separate decision, and one that materially
changes every reported number: the same checkpoint read out and decoded in two
different ways yields different accuracies, a different full-path accuracy and,
in particular, a different measured inconsistency.

Following \S\ref{ssec:problem_train_vs_inference}, the decision factorises into
two independent choices, and the section is organised accordingly:

\begin{enumerate}
    \item a \textbf{readout} $u\mapsto(z_1,\dots,z_L)$, which reduces the
    network's internal taxonomy-node coordinates to one score vector per level
    (\S\ref{ssec:inference_readout});
    \item a \textbf{decoder} $(z_1,\dots,z_L)\mapsto\hat y$, which turns level
    scores into a path (\S\ref{ssec:inference_decoding}).
\end{enumerate}

Both are properties of the evaluation, not of the trained model. Neither touches
the parameters, so both can be replaced after training on a frozen checkpoint ---
which is what makes the grid of \S\ref{ssec:inference_grid} an instrument rather
than a convention. The section then defines the metrics computed on top of a
chosen cell of that grid (\S\ref{ssec:inference_metrics}), states the
checkpoint-selection rule that keeps decoder comparisons honest
(\S\ref{ssec:inference_selection}), and closes by recording exactly which of the
training-time mechanisms can and cannot move which reported quantity
(\S\ref{ssec:inference_interaction}).

\subsection{Readout rules}
\label{ssec:inference_readout}

Every family considered here can be viewed as emitting a vector of
\emph{node coordinates} $u\in\mathbb{R}^{N}$ over the taxonomy, partitioned into
per-level blocks. For the classifier-head families (H-CAST, HRN, HT-CapsNet,
LH-DNN) the blocks are simply the per-level logits of
\S\ref{ssec:problem_scores}; for Hier-COS they are the coordinates of the frame
output. Two readouts are defined on this common representation.

\paragraph{Node score.} Each node is ranked by its own coordinate,
\[
z_l[c] \;=\; \phi\big(u[c]\big),
\]
with $\phi$ the identity for signed classifier logits and $\phi=\lvert\cdot\rvert$
for Hier-COS frame coordinates. The distinction is not cosmetic. A classifier
logit is signed and a large negative value carries the meaning \emph{not this
class}; Hier-COS, by contrast, trains its node distribution as
$\operatorname{log\,softmax}(\lvert u\rvert)$, so its objective never constrains
the sign of a coordinate and only the magnitude carries class evidence. In both
cases the rule ranks nodes by the quantity that the model's own objective drives
up for the correct node.

\paragraph{Subspace norm.} Each node is ranked by the $L_2$ norm of the projection
of $u$ onto the taxonomy subspace $S_c$ assigned to that node, comprising its
ancestors, itself and its descendants:
\[
z_l[c] \;=\; \big\lVert \Pi_{S_c} u \big\rVert_2 .
\]
This is the native rule of Hier-COS (\S\ref{ssec:problem_scores}). Applied to a
classifier-head model it constitutes an \emph{inference-only assumption}: the
per-level logits are treated as coordinates in an identity taxonomy frame. The
assumption is worth stating explicitly because it is not innocuous --- a norm
squares its inputs, so the subspace-norm readout discards the sign of a signed
classifier logit unconditionally. That loss of the ``not this class'' evidence is
the substantive content of the identity-frame assumption, and it is the reason
this readout is expected to behave differently on the two kinds of family.

\paragraph{Each model's native rule is one readout, not a separate case.} The
readouts above are defined for every family, and each family's own inference
procedure is recovered as one of them: \texttt{node\_score} for the
classifier-head families, \texttt{subspace\_norm} for native Hier-COS.
Reproduction is exact for H-CAST and Hier-COS, which recompute the same tensor,
and metric-exact for HRN and the remaining classifier heads, where the readout
reads raw logits rather than the sigmoid or softmax the model applies: those maps
are monotone within a level, and every metric of \S\ref{ssec:inference_metrics} is
computed from an argmax, so no reported quantity changes. Treating the native rule
as a grid cell rather than a privileged baseline is what allows a model to be
scored under another family's readout without special-casing.

\subsection{Decoding rules}
\label{ssec:inference_decoding}

\paragraph{Independent decoding.} Each level is decoded on its own,
\[
\hat y_l=\arg\max_{c}\; z_l[c],\qquad l=1,\dots,L .
\]
No structural information is used. The resulting path need not be consistent: the
predicted fine class may belong to a different parent than the predicted coarse
class. Independent decoding therefore measures what the model has learned at each
level in isolation, and its inconsistency rate is a genuine diagnostic of whether
the training-time mechanisms produced hierarchically coherent scores.

\paragraph{Top-down decoding.} The coarsest level is decoded first and each
subsequent level is restricted to the children of the level above:
\[
\hat y_1=\arg\max_{c}\;z_1[c],
\qquad
\hat y_{l}=\arg\max_{c\,:\,\pi_{l-1}(c)=\hat y_{l-1}}\; z_{l}[c],
\quad l=2,\dots,L .
\]
In the implementation the restriction is a mask setting the scores of
non-children to $-\infty$ before the argmax. Should a predicted parent have no
children in the taxonomy mapping --- which the structural assumption of
\S\ref{sec:methodology_problem} excludes, but which the implementation guards
against --- the level falls back to an unrestricted argmax for the affected
samples.

\begin{proposition}
\label{prop:topdown_consistency}
If the taxonomy satisfies the ``no childless parent'' assumption, the path
produced by top-down decoding is consistent for every input.
\end{proposition}
\begin{proof}
By induction. $\hat y_1$ is unconstrained. For $l>1$, the feasible set
$\{c:\pi_{l-1}(c)=\hat y_{l-1}\}$ is non-empty by assumption, so the masked argmax
is attained on it and $\pi_{l-1}(\hat y_l)=\hat y_{l-1}$ holds by construction; the
fallback branch is therefore never taken.
\end{proof}

The two decoders are not ordered by quality, and the trade-off between them is one
of the objects of study of this thesis. Top-down decoding purchases consistency at
the price of \emph{error propagation}: a mistake at a coarse level removes the
correct fine class from the feasible set, so that a sample which independent
decoding would have classified correctly at the fine level is necessarily
misclassified. Independent decoding avoids this coupling but offers no structural
guarantee. Consequently top-down decoding cannot change the level-$1$ prediction,
can only alter levels $l\ge2$, and may lower fine-level accuracy relative to
independent decoding while lowering the inconsistency rate to zero.

\subsection{The inference grid and its post-hoc use}
\label{ssec:inference_grid}

The readout of \S\ref{ssec:inference_readout} may be preceded by a
\emph{transform} applied to the node coordinates. One transform is considered: the
HCC affine hierarchy projection of \S\ref{sec:methodology_hcc}, applied to $u$
before the readout consumes it. Crossing the two readouts with the presence or
absence of that transform gives four inference rules,
\[
\big\{\texttt{node\_score},\ \texttt{subspace\_norm}\big\}
\times
\big\{\text{none},\ \text{HCC}\big\},
\]
and crossing those with the two decoders of \S\ref{ssec:inference_decoding} gives
the eight cells under which a checkpoint can be scored. Every cell is defined for
every family --- no combination is rejected --- and all of them are computed from
a single shared forward pass, since they differ only in what is done to $u$
afterwards.

\paragraph{Why this is an instrument and not a convention.} Nothing in the grid
trains or modifies a parameter: it is applied to an already-selected checkpoint,
with no loss, optimiser or training loop constructed. That is precisely what makes
it useful. A training-time mechanism that improves the reported metrics admits two
readings --- it changed the representation, or it supplied a better rule for
reading an unchanged representation --- and the grid separates them by
construction. Applying the HCC transform post-hoc to a checkpoint trained
\emph{without} HCC isolates what the constraint is worth as a pure output
correction; comparing that against a checkpoint trained \emph{with} HCC isolates
what the constraint is worth as a training signal. The difference between the two
is the part of HCC's effect that is genuinely representational, and it cannot be
obtained from either run alone. The same logic applies to the frame question of
\S\ref{ssec:design_substrate}: whether Hier-COS's advantage lives in its
subspace-norm readout can be tested by giving that readout to a model that was
never trained with one.

Two properties of the grid follow from the rules themselves, and should not be
mistaken for empirical findings when the results are read.

\begin{proposition}
\label{prop:posthoc_hcc_topdown}
Let the readout be \texttt{node\_score} with $\phi$ the identity, and let the
transform be the HCC projection. Then the transform cannot change any metric
computed under top-down decoding.
\end{proposition}
\begin{proof}
By \S\ref{ssec:hardnet_to_hcc} the correction is
$\hat z_{l+1} = z_{l+1} - M^{\top}(MM^{\top}+\varepsilon I)^{-1}(Mz_{l+1}-\hat z_l)$.
By Lemma~\ref{lemma:diag} the inverted matrix is diagonal, so the vector
subtracted from $z_{l+1}$ assigns to every child the same value, namely the
coefficient of its parent; within the children of a fixed parent the correction is
therefore a common additive shift, which preserves their ordering. Top-down
decoding evaluates an argmax restricted to exactly such a sibling set
(\S\ref{ssec:inference_decoding}). The coarsest level is anchored, $\hat z_1=z_1$,
so $\hat y_1$ is unchanged; inductively, each subsequent masked argmax is taken
over a set whose internal ordering the correction preserves, so $\hat y_l$ is
unchanged for every $l$.
\end{proof}

The proposition fails for the magnitude convention $\phi=\lvert\cdot\rvert$ used
by Hier-COS, where a common additive shift is not order-preserving, and it fails
for the subspace-norm readout, whose subspaces span several levels: there the
transform moves both decoders for every family. The practical consequence is a
check rather than a result --- an \texttt{hcc\_node\_score} row differing from its
untransformed counterpart under top-down decoding on a classifier-head model
indicates a defect in the evaluation, not an effect.

\paragraph{Three qualifications on the post-hoc HCC cells.} First, the transform
is always applied at $\alpha=1$ and does not replicate a run's activation
schedule; for a checkpoint trained with HCC the schedule-faithful reference is the
alpha its own forward pass applied at test time, recorded as
\texttt{proj\_constraint\_alpha}, and the post-hoc cell is compared against that
rather than assumed to reproduce it. Second, the transform acts on whatever
coordinates the readout consumes, which for HRN is its auxiliary leaf
cross-entropy head rather than the sigmoid tree branch that HCC corrects during
training; the HRN cells are therefore not a reconstruction of that family's
training-time HCC path and are not reported as one. Third, the grid inherits the
hazards of any post-hoc tool applied to historical runs: where a run's recorded
test manifest no longer resolves and the official dataset adapter is substituted,
the rebuilt label space may differ from the one the checkpoint trained on, which
surfaces as plausible but meaningless parent-level metrics while leaf accuracy
still matches. Every reported cell is therefore checked against the run's own
recorded test metrics on its native cell before any other row from that run is
used.

\subsection{Metrics}
\label{ssec:inference_metrics}

The metrics below are functions of a decoded path, so they are defined relative to
a chosen cell of the grid of \S\ref{ssec:inference_grid}. Unless a table states
otherwise, results are reported under each family's native readout, and always
under both decoders as separate quantities; rows from other cells of the grid are
labelled as such. Writing $\hat y=(\hat y_1,\dots,\hat y_L)$ for the decoded path and
$y$ for the ground-truth path:

\begin{itemize}
    \item \textbf{Per-level accuracy.} $\operatorname{Acc}_l
    =\Pr[\hat y_l=y_l]$, reported for each level. \emph{Higher is better.}
    \item \textbf{Weighted average precision (wAP).} The mean of the per-level
    accuracies weighted by the number of classes at each level,
    \[
    \mathrm{wAP}=\frac{\sum_{l}C_l\,\operatorname{Acc}_l}{\sum_l C_l}.
    \]
    The weights are the class counts $C_l$ and are unrelated to the objective's
    level weights $w_l$ of \S\ref{ssec:problem_objective}; since $C_L\gg C_1$ in
    all datasets used here, wAP is dominated by the finest level.
    \emph{Higher is better.}
    \item \textbf{Full-path accuracy (FPA).} The exact-match rate over the whole
    path, $\mathrm{FPA}=\Pr[\hat y_l=y_l\ \text{for all }l]$. It is the strictest
    of the accuracy measures and the primary quantity for checkpoint selection.
    \emph{Higher is better.}
    \item \textbf{Average hierarchical distance (AHD).} The expected number of
    levels separating the prediction from the ground truth, measured as the depth
    below the deepest correctly predicted ancestor,
    \[
    \mathrm{AHD}=\mathbb{E}\Big[\,L-\max\{m: \hat y_l=y_l\ \text{for all } l\le m\}\Big],
    \]
    that is, $L$ minus the length of the longest matching prefix. It equals $0$ on
    a fully correct path and $L$ when even the coarsest level is wrong, and it
    distinguishes a near miss within the correct coarse class from a mistake at the
    top of the taxonomy. \emph{Lower is better.}
    \item \textbf{Tree-induced consistency error (TICE).} The fraction of samples
    whose decoded path is \emph{not} consistent,
    $\mathrm{TICE}=1-\Pr[\pi_l(\hat y_{l+1})=\hat y_l\ \text{for all }l]$. It is a
    property of the prediction alone and does not involve the labels.
    \emph{Lower is better.}
\end{itemize}

By Proposition~\ref{prop:topdown_consistency}, $\mathrm{TICE}=0$ identically under
top-down decoding. The informative quantity is therefore
$\mathrm{TICE}_{\text{independent}}$, which measures whether the model produces
coherent scores on its own; a comparison of TICE across methods is meaningful only
in the independent column, and reporting the top-down column serves as a check
that the decoder is behaving as specified.

\subsection{Checkpoint selection}
\label{ssec:inference_selection}

Because the two decoders induce different orderings of the epochs, a single
``best'' checkpoint would silently favour whichever decoder the selection metric
happened to reward. Every run therefore maintains \emph{two} checkpoints, one per
decoding mode, each selected on the validation split by the mode's own metrics.

The selection criterion is itself lexicographic, which is convenient given the
subject of this thesis. Epochs are ranked by the tuple
\[
\Big(\;\mathrm{FPA}_{\text{mode}},\;\; -\mathrm{TICE}_{\text{mode}},\;\;
\mathrm{wAP}_{\text{mode}}\;\Big)
\]
under the usual lexicographic order on tuples, so that full-path accuracy
dominates, consistency breaks ties, and weighted accuracy breaks any remaining
tie. Missing or non-finite entries rank below every finite value; when none of the
three hierarchy metrics is available the deepest-level accuracy is used as the
sole criterion.

The reporting rule that follows is applied without exception in the experimental
chapters: \textbf{a top-down result is read from the top-down-selected checkpoint
and an independent result from the independent-selected checkpoint}. The two rows
of a table may therefore come from different epochs of the same run, and mixing
them --- reporting, for instance, the independent-decoded metrics of the
top-down-selected checkpoint --- would confound the decoder comparison with a
checkpoint comparison.

\subsection{Interaction with the training-time mechanisms}
\label{ssec:inference_interaction}

Combining \S\ref{ssec:problem_train_vs_inference} with the results established
above yields a clean division, which is worth stating because it is easy to
assume otherwise.

\begin{itemize}
    \item \textbf{Lexicographic mode and the LH-projection have no effect at
    inference.} The first modifies the update direction by construction; the
    second does so by Proposition~\ref{prop:lhproj_forward}. Both change the
    parameters reached by training, and neither changes the function computed by a
    given parameter vector. Any difference they produce in the metrics above is
    entirely mediated by the learned parameters.
    \item \textbf{HCC does act at inference}, since it modifies the logits
    themselves and the projection is applied in evaluation as in training, with
    $\alpha_t$ at its final value.
    \item \textbf{HCC does not by itself guarantee consistency.} This is the point
    most easily misread. The cascade enforces $M_{l,l+1}\hat z_{l+1}=\hat z_l$, an
    equality in \emph{logit space} between the parent logit and the sum of its
    children's logits, as emphasised in \S\ref{ssec:hardnet_to_hcc}. It does not
    imply
    $\pi_l\big(\arg\max_j \hat z_{l+1}[j]\big)=\arg\max_i \hat z_l[i]$: a parent
    may hold the largest logit while the largest child logit overall belongs to a
    different parent whose children sum to less. HCC should therefore be expected
    to \emph{reduce} $\mathrm{TICE}_{\text{independent}}$, not to zero it, and any
    claim that it produces consistent predictions must be supported by the measured
    independent-decoding TICE rather than deduced from the constraint. Exact
    consistency at inference is delivered by top-down decoding alone.
\end{itemize}

The division has a methodological consequence that
\S\ref{sec:methodology_design} relies on. Because lexicographic mode and the
LH-projection cannot act at inference, any effect they produce is necessarily
representational, and no post-hoc analysis can either reproduce or explain it
away. Because HCC \emph{can} act at inference, the same is not true of it: an
improvement under HCC is a priori compatible with a purely corrective reading,
and separating the two requires the post-hoc grid of
\S\ref{ssec:inference_grid}. The asymmetry is not a nuisance but the reason the
three mechanisms cannot be compared on their headline numbers alone.



\section{Experimental design}
\label{sec:methodology_design}

The chapter has so far introduced three mechanisms for imposing a coarse-to-fine
priority, each at a different locus, and the inference apparatus under which any
of them is scored. What remains is to turn that apparatus into an experiment.

This section does so in the order in which the constraints actually bind. It
first collects the three mechanisms into a single comparison and establishes that
they form a mutually exclusive choice rather than a factorial
(\S\ref{ssec:design_mechanisms}--\S\ref{ssec:design_axis}); it records which
families admit which mechanism (\S\ref{ssec:design_support}); it shows how the
nominal design space is cut down by structural constraints, several of which are
themselves findings (\S\ref{ssec:design_space}); it states the research questions
that the surviving configurations answer, grouped by what they are questions
\emph{about} (\S\ref{ssec:design_substrate}--\S\ref{ssec:design_locus}); it
isolates the arms that exist only to falsify an attribution
(\S\ref{ssec:design_controls}); and it closes with the variables held fixed, the
reporting protocol, and the threats to validity
(\S\ref{ssec:design_threats}--\S\ref{ssec:design_protocol}).

\subsection{Where each mechanism acts}
\label{ssec:design_mechanisms}

\begin{table}[htbp]
\centering
\small
\begin{tabular}{@{}lllll@{}}
\toprule
\textbf{Intervention} & \textbf{Acts on} & \textbf{Granularity} & \textbf{Guarantee} & \textbf{At inference} \\
\midrule
HCC & logits & per sample & exact (affine equality) & yes \\
LH-projection & update direction & per sample & first order, at the branch point & no \\
Lexicographic mode & update direction & per batch & first order, per block & no \\
\midrule
Readout rule & level scores & per sample & --- & yes, only \\
Top-down decoding & decoded path & per sample & exact (path consistency) & yes, only \\
\bottomrule
\end{tabular}
\caption{The interventions, classified by the point of the computation at which
they act. The upper block contains the three priority mechanisms, the lower block
the two components of the inference rule. Only HCC and the inference rules change
the prediction produced by a fixed parameter vector; the two gradient-space
mechanisms act solely through the parameters that training reaches.}
\label{tab:framework_mechanisms}
\end{table}

The three priority mechanisms form the axis announced in
\S\ref{ssec:lexicographic}: priority imposed on the output contract by
construction (HCC), priority imposed on the learning signal at a structural
bottleneck (the LH-projection), and priority imposed on the assembled update in
parameter space (lexicographic mode). None of them achieves sequential
lexicographic optimality; all three realise lexicographic \emph{structure} at a
different locus. Whether the locus matters is the question the design exists to
answer.

The lower block of the table is what makes that question answerable. Because the
readout and the decoder can be varied on a frozen checkpoint
(\S\ref{ssec:inference_grid}), the contribution of a mechanism to the learned
representation can be distinguished from its contribution to the readout, which
for HCC --- the one mechanism active at inference --- is not otherwise separable.

\subsection{The mechanism axis is a single choice, not a factorial}
\label{ssec:design_axis}

It is natural to expect three independently configurable mechanisms to yield a
$2\times2\times2$ design. They do not. The three are pairwise incompatible, and
because the incompatibilities all follow from one property rather than from
implementation accident, the fact is stated here as a single rule.

\begin{itemize}
    \item \textbf{HCC and the LH-projection} both determine the per-level logits,
    by incompatible routes: the LH-projection fixes what each head reads, while
    HCC overwrites what each head emits with the solution of an affine problem
    anchored on the level above. Under HCC the head weights whose row space the
    LH-projection reserves no longer describe the map that produces the logits the
    objective sees. The pair is rejected at construction, for the reason given in
    \S\ref{ssec:lhproj_hiercos}.
    \item \textbf{The LH-projection and lexicographic mode} are, by the analysis
    of \S\ref{ssec:lhproj_forward}, two implementations of one principle: both
    remove from a lower-priority level's learning signal the component lying in a
    subspace claimed by the higher levels, one inside the computational graph and
    one on the assembled gradient. Composing them would apply the same principle
    twice by two routes, and the resulting ordering would not be attributable to
    either. The pair is likewise rejected at construction.
    \item \textbf{HCC and lexicographic mode} are incompatible for a reason that
    is easy to miss, because the loci argument appears to permit the combination.
    Lexicographic mode presupposes three \emph{separable} level objectives, whose
    gradients can be computed independently and then ordered
    (\S\ref{ssec:lex_setting}). Activating HCC destroys that premise: the level
    losses are evaluated on the effective logits $z_l^{\mathrm{eff}}$, and by the
    cascade of \S\ref{ssec:hardnet_to_hcc} each of those depends on the raw logits
    of every level at or above $l$. There are then no three separable objectives
    for the ordering to act upon, and the priority imposed on gradients that HCC
    has already cross-coupled would not mean what the construction claims. The
    combination is not run in this thesis.
\end{itemize}

The mechanism axis is therefore a four-way single choice,
\[
\big\{\,\text{none},\ \text{HCC},\ \text{LH-projection},\ \text{lexicographic mode}\,\big\},
\]
and every comparison among the three is a between-arms comparison against the
shared \emph{none} baseline.

Two remarks are due. First, this is not a limitation of the design but one of its
results: hierarchical priority can be bought in output space \emph{or} in gradient
space, and the two are not composable, because the output-space constraint
destroys the objective separability that gradient-space ordering requires. That
the three exclusions reduce to a single property --- separability of the level
objectives, violated from three directions --- is the cleanest statement the
thesis can make about how these families of method relate. Second, the exclusion
of HCC with lexicographic mode is a methodological commitment and not, at the time
of writing, an enforced one: the implementation rejects the other two pairs at
construction but accepts this one, so the guarantee rests on the configurations
actually run rather than on a validator. No configuration or run in this thesis
enables both.

\subsection{Which families support which mechanism}
\label{ssec:design_support}

\begin{table}[htbp]
\centering
\small
\begin{tabular}{@{}lccc l@{}}
\toprule
\textbf{Family} & \textbf{HCC} & \textbf{Lex.\ mode} & \textbf{LH-proj.} & \textbf{Condition} \\
\midrule
H-CAST   & yes & yes & no  & three levels; \texttt{globalkl:\ false} for lex.\ mode \\
Hier-COS & yes & yes & yes & decomposed loss for lex.; identity frame for LH-proj. \\
HRN      & yes & yes & no  & exactly three levels; decomposed level-marginal loss \\
HT-CapsNet & yes & yes & no & three capsule margin losses \\
LH-DNN   & no  & no  & built in & no three independent level objectives \\
\bottomrule
\end{tabular}
\caption{Support matrix. ``Built in'' indicates that the mechanism is constitutive
of the family rather than optional.}
\label{tab:framework_support}
\end{table}

Three entries deserve comment, and the first is the most consequential for the
design.

\paragraph{LH-DNN carries no mechanism arms, and this is an argument rather than a
gap.} Every mechanism in this thesis acts on the part of the network where a
shared representation is distributed to level-specific heads --- which is
precisely the part that constitutes LH-DNN's contribution. Enabling lexicographic
mode on it would apply, by explicit gradient surgery, the same principle its
branch-point projection already applies inside the graph; imposing HCC on it would
overwrite the structure it exists to demonstrate. Any arm applied to LH-DNN
therefore either duplicates what it already does or replaces it, and a result from
such an arm could not be attributed to LH-DNN. It is retained as a baseline and as
the origin of the LH-projection, and it is the right control precisely because it
admits no arms: it holds the position ``hierarchy handled natively by the
architecture'' fixed while the other four families vary.

\paragraph{Hier-COS is the only family supporting all three mechanisms}, which is
why it carries the structural ablations of \S\ref{ssec:lhproj_hiercos} and serves
as the common substrate on which the loci can be compared at matched backbone,
frame, loss and weights. The head-to-head between the LH-projection and
lexicographic mode --- the two implementations of one principle --- is therefore
available \emph{within} a single family, holding everything but the implementation
fixed, which is a cleaner contrast than a cross-family comparison could have been.

\paragraph{HRN's three-level restriction} is a property of its implementation, not
of the hierarchy, and coincides here with the three-level restriction that HCC and
lexicographic mode impose for their own reasons.

\subsection{The design space and what prunes it}
\label{ssec:design_space}

Enumerated as a factorial, the space is intractable and uninformative. On Hier-COS
alone the frame, the loss scope, the level weighting, the learnable transform, the
mechanism choice and the dataset already give
$3\times3\times3\times3\times7\times3 = 1701$ configurations per seed, before any
backbone or head variant, and at three seeds each. The factorial is not the plan.

What cuts it to the order of a few dozen configurations that carry information is
a set of structural constraints, enforced as validation errors rather than
observed as conventions. They are worth stating in the thesis because several of
them are themselves substantive:

\begin{enumerate}
    \item \textbf{Gradient-space ordering requires separable level objectives.}
    Lexicographic mode requires exactly three differentiable level losses, so a
    loss with a single coupled objective over all taxonomy nodes cannot be used
    with it. The same requirement, seen from the other side, forces the
    LH-projection to a per-level softmax: a global softmax over all $N$ nodes
    couples every node into one normaliser, and there is then no separable
    per-level gradient for a branch head to receive. Constraints on lexicographic
    mode and on the LH-projection are one fact viewed from two directions, and it
    is the same fact that excludes HCC from the combination
    (\S\ref{ssec:design_axis}).
    \item \textbf{The LH-projection requires independent level heads}, hence an
    identity or per-level block-diagonal frame, and a representation wider than
    the total number of non-leaf classes (Proposition~\ref{prop:lhproj_rank}).
    \item \textbf{HCC requires exactly three levels}, as do the three-objective
    constructions of lexicographic mode.
    \item \textbf{Family-specific decompositions}, recorded in
    Table~\ref{tab:framework_support}: \texttt{globalkl:\ false} on H-CAST, the
    level-marginal objective on HRN, the three capsule margin losses on
    HT-CapsNet.
\end{enumerate}

The pruning is therefore not a convenience. The constraints trace the boundary of
where each mechanism is definable at all, and \S\ref{ssec:design_substrate}
treats that boundary as an object of study rather than as a precondition to be
satisfied silently.

\subsection{Substrate questions: what the comparison runs on}
\label{ssec:design_substrate}

Before the mechanisms can be compared, the platform they are compared on has to be
characterised. Three questions do that, and each has a second reading as an
enabling condition, which is what makes them part of the argument rather than
preliminaries.

\paragraph{Is Hier-COS's benefit the fixed frame, or the inference rule?} Hier-COS
bundles three things that are usually credited as one: a frozen orthonormal frame
over taxonomy nodes, a subspace-norm readout (\S\ref{ssec:inference_readout}), and
a loss with a level regulariser. The published contribution narrative attaches to
the frame. The contrast that tests this is a three-point ladder on the frame,
holding everything else fixed --- a dense random orthonormal frame, a per-level
block-diagonal one, and the identity --- which removes cross-level coordinate
mixing and then the rotation entirely, while leaving the node-block layout and the
readout in place. If the identity matches the dense random frame, the frame is not
doing the work and the honest claim becomes that the gain is attributable to the
output parametrisation and the readout; if it does not, the mechanism is
cross-level mixing, and the block-diagonal rung separates ``orthonormality helps''
from ``mixing helps''. The post-hoc grid sharpens this considerably: because the
subspace-norm readout can be applied to any family
(\S\ref{ssec:inference_readout}), the readout half of the bundle can be tested on
models never trained with a frame at all, without retraining anything.

The second reading: the block-diagonal rung is a \emph{precondition} of the
LH-projection, not an optional extra, so this ladder is a prerequisite for the
locus comparison of \S\ref{ssec:design_locus} rather than a standalone ablation.

\paragraph{Where must the softmax live?} Normalising across all taxonomy nodes at
once, against normalising within each level, looks like a loss detail and is in
fact the structural gate for most of the design (\S\ref{ssec:design_space}). It
carries two questions. As a standalone ablation: does per-level normalisation cost
accuracy relative to global, at matched frame and weights? As an enabling
condition: per-level normalisation is what makes per-level weighting, branch
heads, and gradient-space ordering definable at all --- so if it costs accuracy,
that cost is the \emph{entry price} of every gradient-space method in this thesis
and must be reported as such, not absorbed silently into a
mechanism-versus-baseline delta. This is the reason the baseline is run in both
scopes: a mechanism arm at per-level softmax compared against a baseline at global
softmax confounds the mechanism with the normaliser.

\paragraph{Does gradient-space ordering need a learnable transform to act on?}
Progressively removing the learnable transformation between backbone and frame
--- from the full residual transform, to a batch-norm-plus-linear reduction, to
none at all --- is a mechanism question, not a capacity question. Gradient
projection acts on whatever parameters lie in the projected path. If the benefit
shrinks as the transform shrinks, the effect is localised to a learned
re-embedding; if it survives with no transform in the optimiser path, the
mechanism is acting on the backbone itself, and the claim generalises to any
architecture --- which is what the cross-family extension of
Table~\ref{tab:framework_support} depends on.

\subsection{Locus questions: does it matter where priority is imposed?}
\label{ssec:design_locus}

These are the questions the chapter was built to ask, and they map one-to-one onto
the loci of \S\ref{ssec:problem_train_vs_inference}.

\paragraph{Does gradient-space ordering work, and is the \emph{ordering} what
matters?} The first half is a comparison against the matched baseline. The second
half is the harder one and requires the composition modes of
\S\ref{ssec:lex_modes} to be treated as an experimental axis rather than an
implementation choice: \texttt{coarse\_first} is the hypothesis,
\texttt{fine\_first} inverts the priority, and \texttt{pairwise\_orthogonal}
decorrelates the level gradients while imposing the weakest ordering. Without the
last two, an observed improvement is equally compatible with ``projecting
gradients at all helps''. A separate axis asks whether the projection needs to be
unconditional: gating it to fire only when the level gradients actually conflict
would, if it matched unconditional projection, license a far more mechanistic
claim --- that the benefit is conflict resolution and the operator is inert most
of the time --- and that claim is checkable directly against the logged gate
statistics of \S\ref{ssec:lex_implementation}. Finally, the activation epoch of
\S\ref{ssec:lex_activation} asks whether ordering is a training-dynamics
intervention or a late correction, and deliberately mirrors the HCC schedule so
that the two families remain comparable on that axis.

\paragraph{One principle, two implementations: LH-projection against lexicographic
mode.} By \S\ref{ssec:lhproj_forward} these are not complementary methods but the
same idea realised twice. The in-graph formulation is elegant and cheap --- the
constraint is expressed once in the forward pass and the backward pass inherits it
--- but the guarantee it obtains is only what the graph gives: first-order, per
sample, at the branch point, and conditional on the activation. The explicit
formulation ignores how the graph is built, takes whatever gradients emerge and
orthogonalises them directly; it costs $L$ backward passes and three simultaneous
gradient copies, and buys an ordering that holds over the whole objective set with
no activation assumption. The question is therefore not whether they stack --- they
cannot (\S\ref{ssec:design_axis}) --- but whether the cheap in-graph formulation
buys the same ordering, or whether its branch-point restriction costs real
hierarchy consistency. That is a head-to-head at matched frame, loss, weights and
seeds, scored primarily on TICE and AHD against the shared baseline, and the
sub-variants of \S\ref{ssec:lhproj_hiercos} probe exactly where its guarantee is
thin: the sample-dependent reserved subspace tests whether the activation
assumption bites, and the representation width sets how much room exists above the
reserved subspace.

\paragraph{Output space against gradient space.} This is the contrast that keeps
the thesis honest about what ``lexicographic'' means. HCC constrains outputs; the
other two constrain gradients. Per \S\ref{ssec:design_axis} the achievable cells
are exactly the four of the single-choice axis, so this is a clean between-arms
comparison against a shared baseline, and the absence of combination arms is
itself the result. Two sub-questions refine it. The activation schedule asks
whether the output constraint should be imposed from initialisation or applied as
a late correction. More decisively, the constraint can be applied \emph{only} at
evaluation --- either by a run configured for test-time-only enforcement, or, at
no training cost at all, by the post-hoc HCC cells of \S\ref{ssec:inference_grid}
applied to a checkpoint trained without it. That comparison separates ``the
constraint improved the learned representation'' from ``the constraint fixed up
the predictions'', which is the single most important distinction to establish
before any claim about HCC is made, and Proposition~\ref{prop:posthoc_hcc_topdown}
fixes in advance which decoder can register such an effect.

\subsection{Controls and falsification arms}
\label{ssec:design_controls}

Three configurations exist solely to falsify an attribution, and are recorded here
so that they are not mistaken for hypotheses.

\begin{enumerate}
    \item \textbf{Inverted priority.} The \texttt{fine\_first} composition mode of
    \S\ref{ssec:lex_modes} applies exactly the same quantity of gradient surgery
    with the ordering reversed. A gain reproduced under it cannot be attributed to
    the hierarchy, only to orthogonalisation as such.
    \item \textbf{Order-free projection.} \texttt{pairwise\_orthogonal}
    decorrelates the level gradients under the weakest ordering the operator
    admits. It is what separates ``coarse-first ordering helps'' from ``projecting
    at all helps'', and without it the central claim of
    \S\ref{ssec:design_locus} is not identified.
    \item \textbf{Coarse-heavy weighting.} If the reported benefit of
    gradient-space ordering is consistency rather than accuracy, then reweighting
    the scalarised objective towards the coarse levels is a much cheaper way to buy
    the same thing. A coarse-heavy weight vector tests exactly that. If it
    reproduces the consistency gain without any projection, the claim about
    gradient-space ordering weakens sharply --- it would mean the mechanism is one
    way to move along an accuracy--consistency trade-off and simple reweighting is
    another.
\end{enumerate}

\subsection{Controlled variables and threats to validity}
\label{ssec:design_threats}

\paragraph{Level weights.} The weighting $w_l$ of \S\ref{ssec:problem_objective}
is the most consequential controlled variable, for a reason internal to the
design: activating lexicographic mode replaces the configured weighting by a
uniform one, because the tensors consumed by the projection are the raw level
objectives (\S\ref{ssec:lex_setting}). A lexicographic arm and a baseline arm that
differ in their configured weighting therefore differ in two respects at once, and
the comparison isolates ``priority plus uniform weighting'' against ``weighted
scalarisation'', not priority alone. Where such pairs occur they are identified
explicitly, and matched-weight comparisons are reported wherever available. Two
questions must be kept apart here: whether the weight choice is a large effect,
and whether the \emph{ranking of methods} is weight-dependent. The second is the
substantive one --- a method that wins only at one weight setting has not been
shown to win.

\paragraph{Backbone and pretraining.} A from-scratch, CIFAR-native wide residual
backbone and an ImageNet-pretrained ResNet-50 bound every other conclusion. This
axis is not a contribution; it exists to establish whether the frame, weighting
and mechanism effects are conditional on the feature extractor. If they invert
between the two, every conclusion in the thesis is conditional in the same way.

\paragraph{Dataset scope and local extrapolations.} The three datasets differ in
how much the hierarchy actually constrains the fine label: a coarse semantic
grouping over small images, a fine-grained set with shallow semantic spread, and a
genuinely hierarchical manufacturer--family--variant taxonomy. That is the property
every mechanism here depends on, so a method that works on one and not another is
reporting a requirement, not a failure. Separately, several settings used here
depart from the configurations the original works evaluated --- among them HRN on
CIFAR-100, Hier-COS on CUB, and this repository's construction of the CIFAR
hierarchy. Results from those settings are reported separately from paper-aligned
ones and are not used to support claims about the published methods.

\paragraph{Coverage.} The design above describes the space the apparatus admits,
not a claim that every cell has been populated. Where an arm required by a
question has not been run, the corresponding claim is reported as unidentified
rather than as unsupported, and the experimental chapter states which cells are
missing and what each absence costs.

\subsection{Protocol and reporting rules}
\label{ssec:design_protocol}

The following are applied without exception, and several of them exist to prevent
a specific error that the apparatus of this chapter makes possible.

\begin{itemize}
    \item \textbf{One mechanism at a time.} The main experiments vary a single
    mechanism against a matched baseline, so that an effect can be attributed to a
    locus. By \S\ref{ssec:design_axis} no combination arm exists among the three
    priority mechanisms, and none is planned.
    \item \textbf{Both decoders are always reported}, for every configuration,
    each from its own checkpoint (\S\ref{ssec:inference_selection}): a top-down row
    from the top-down-selected checkpoint, an independent row from the
    independent-selected one. Mixing them would confound the decoder comparison
    with a checkpoint comparison. A decoder is never treated as a property of a
    method.
    \item \textbf{The structural preconditions of the LH-projection are not
    ablations of it.} The identity frame and the decomposed level objective are
    what make the mechanism definable on Hier-COS at all
    (\S\ref{ssec:lhproj_hiercos}). Their comparison against the published
    configuration measures the cost of adapting the baseline, and is reported
    before, and separately from, any comparison involving the projection itself.
    \item \textbf{Mechanism activation is verified from diagnostics, not from
    configuration.} That a run was configured with a mechanism does not establish
    that the mechanism was active: HCC activation is read from the logged
    $\alpha_t$, and claims about gradient-space ordering are grounded in the logged
    cosines, norms and guard flags of \S\ref{ssec:lex_implementation} rather than
    in a configuration flag. This matters most where a schedule delays activation,
    since a directory name or a preset name records the intent and not the outcome.
    \item \textbf{Metric directionality is fixed.} Full-path accuracy, weighted
    average precision and per-level accuracy are higher-is-better; average
    hierarchical distance and tree-induced consistency error are lower-is-better.
    Differences between percentages are reported in percentage points.
    \item \textbf{Seeds.} Configurations are run at three seeds unless otherwise
    stated, and aggregate quantities are reported with the seed count and the
    sample standard deviation alongside the mean.
    Epoch-level curves are filtered for runs whose logs terminate early before
    being averaged, since a truncated log is not a truncated run.
\end{itemize}
